\documentclass[11pt,oneside]{article}

\usepackage[a4paper,margin=27mm,headheight=15pt]{geometry}
\usepackage[T1]{fontenc}
\usepackage[utf8]{inputenc}
\IfFileExists{libertinus.sty}{\usepackage{libertinus}}{\usepackage{lmodern}}
\usepackage{microtype}
\usepackage{amsmath,amssymb,amsthm,mathtools,mathrsfs}
\usepackage{bm}
\usepackage{booktabs,longtable,array,tabularx}
\usepackage{graphicx}
\usepackage{xcolor}
\usepackage{enumitem}
\usepackage{fancyhdr}
\usepackage{titlesec}
\usepackage{listings}
\usepackage{xurl}
\usepackage{hyperref}
\usepackage[nameinlink,noabbrev]{cleveref}

\definecolor{linkblue}{RGB}{25,75,135}
\definecolor{shadegray}{RGB}{246,247,249}
\definecolor{rulegray}{RGB}{195,201,208}
\hypersetup{
  colorlinks=true,
  linkcolor=linkblue,
  citecolor=linkblue,
  urlcolor=linkblue,
  pdftitle={A Formula Atlas for the Thue--Morse Sequence},
  pdfauthor={Vladimir Reshetnikov},
  pdfsubject={Digital, valuative, Pascal, generating, Prouhet, Fourier, spectral, and arithmetic representations},
  pdfkeywords={Thue-Morse, Prouhet, binary digit sum, Lucas theorem, Kummer theorem, Mahler function, Walsh transform, Riesz product, Fabius function}
}

\pagestyle{fancy}
\fancyhf{}
\fancyhead[L]{\small Thue--Morse formula atlas}
\fancyhead[R]{\small\nouppercase{\leftmark}}
\fancyfoot[C]{\thepage}
\renewcommand{\headrulewidth}{0.4pt}

\titleformat{\section}{\Large\bfseries}{\thesection}{0.7em}{}
\titleformat{\subsection}{\large\bfseries}{\thesubsection}{0.7em}{}
\titleformat{\subsubsection}{\normalsize\bfseries}{\thesubsubsection}{0.7em}{}
\setcounter{secnumdepth}{3}
\setcounter{tocdepth}{2}
\setlist[itemize]{leftmargin=2em,itemsep=0.25em,topsep=0.4em}
\setlist[enumerate]{leftmargin=2.3em,itemsep=0.25em,topsep=0.4em}
\setlength{\emergencystretch}{2em}
\allowdisplaybreaks

\newtheorem{theorem}{Theorem}[section]
\newtheorem{proposition}[theorem]{Proposition}
\newtheorem{lemma}[theorem]{Lemma}
\newtheorem{corollary}[theorem]{Corollary}
\newtheorem{identity}[theorem]{Identity}
\theoremstyle{definition}
\newtheorem{definition}[theorem]{Definition}
\newtheorem{algorithm}[theorem]{Algorithm}
\newtheorem{example}[theorem]{Example}
\theoremstyle{remark}
\newtheorem{remark}[theorem]{Remark}
\newtheorem{warning}[theorem]{Warning}

\newcommand{\R}{\mathbb R}
\newcommand{\C}{\mathbb C}
\newcommand{\Q}{\mathbb Q}
\newcommand{\N}{\mathbb N}
\newcommand{\Z}{\mathbb Z}
\newcommand{\Ftwo}{\mathbb F_2}
\newcommand{\e}{\mathrm e}
\newcommand{\ii}{\mathrm i}
\newcommand{\dd}{\,\mathrm d}
\newcommand{\wt}{s_2}
\newcommand{\vtwo}{v_2}
\newcommand{\eps}{\varepsilon}
\newcommand{\xor}{\mathbin{\oplus}}
\newcommand{\bitand}{\mathbin{\&}}
\newcommand{\submask}{\mathrel{\preccurlyeq_2}}
\newcommand{\ind}{\mathbf 1}
\newcommand{\sgn}{\operatorname{sgn}}
\newcommand{\res}[2]{\left[#1\right]_{#2}}
\newcommand{\Bell}{\mathcal B}
\newcommand{\bigO}{\mathcal O}
\newcommand{\smallo}{o}
\newcommand{\Up}{\operatorname{up}}
\newcommand{\repo}{\nolinkurl{Analysis/FabiusFunction}}
\newcommand{\lean}[1]{%
  \ifmmode\text{\nolinkurl{#1}}\else\nolinkurl{#1}\fi}

\newenvironment{proofidea}{\par\smallskip\noindent\textit{Proof architecture.}\ }{\hfill$\square$\par\smallskip}
\newenvironment{boxedremark}{%
  \par\smallskip\noindent\begin{center}\begin{minipage}{0.94\textwidth}%
  \color{black}\hrule\vspace{0.6em}\small
}{%
  \vspace{0.6em}\hrule\end{minipage}\end{center}\smallskip
}

\lstdefinestyle{pseudo}{
  basicstyle=\ttfamily\small,
  backgroundcolor=\color{shadegray},
  frame=single,
  rulecolor=\color{rulegray},
  columns=fullflexible,
  keepspaces=true,
  showstringspaces=false,
  xleftmargin=0.7em,
  xrightmargin=0.7em,
  aboveskip=0.8em,
  belowskip=0.8em
}

\title{\bfseries A Formula Atlas for the Thue--Morse Sequence\\[0.35em]
\large Digital, valuative, Pascal, generating, Prouhet, Fourier, spectral, and arithmetic representations}
\author{Consolidated companion report to the Fabius--Rvachev formal development\\
\small Vladimir Reshetnikov's \texttt{ProveIt} repository}
\date{26 August 2026}

\begin{document}
\maketitle

\begin{abstract}
Four independent drafts in the \texttt{ProveIt} Fabius-function research-frontier directory
studied exact formulae for the Thue--Morse sequence from overlapping but nonidentical
viewpoints.  This article merges them into one coherent exposition, removes duplicated
proofs, reconciles notation and normalizations, and retains the strongest version of each
result.  It develops the signed sequence
\(\eps_n=(-1)^{\wt(n)}\) and the zero--one bit
\(\tau(n)=\wt(n)\bmod2\) through binary recurrences, floor and fractional-part sums,
trigonometric signs, factorial and binomial valuations, carry cocycles, Pascal support,
multinomial parity, Boolean M\"obius inversion, odious and evil enumerators, lacunary Euler
products, nonlinear recurrences, Bell polynomials, Hessenberg determinants, algebraic power
series over \(\Ftwo\), finite differences, exact Prouhet moments, sparse Pascal-row
cancellation, ordinary and Walsh Fourier transforms, Riesz products, autocorrelation,
Dirichlet transforms, and evaluated infinite products.

Several additions go beyond all four drafts.  We derive exact dyadic formulas for
Thue--Morse sums in residue classes modulo three and place them inside Newman's phenomenon;
record the overlap-free and substitution-dynamical structure of the infinite word; explain
the nonvanishing Hankel-determinant theorem and its Pad\'e consequences; give a direct
natural-boundary argument for the Mahler product; and collect the transcendence and sharp
irrationality-exponent results for Thue--Morse--Mahler numbers.  A final section maps the
paper results to existing Lean modules and proposes a dependency-aware formalization order
for the additional material.

Every elementary identity is proved in the text.  Results whose known proofs require deeper
symbolic dynamics, Mahler's method, or determinant theory are explicitly identified as
literature theorems.  No statement is described as formally verified unless a corresponding
Lean declaration or module is named.
\end{abstract}

\tableofcontents
\newpage

\section{Scope, conventions, and the formula being complemented}
\label{sec:scope}

Throughout, \(\N=\{0,1,2,\ldots\}\).  Write the binary expansion of \(n\) as
\begin{equation}
  n=\sum_{j\ge0}b_j(n)2^j,
  \qquad b_j(n)\in\{0,1\},
  \qquad \wt(n)=\sum_{j\ge0}b_j(n).                         \label{eq:binary-expansion}
\end{equation}
Only finitely many digits are nonzero.  The two standard normalizations of the
Thue--Morse sequence are
\begin{equation}
  \tau(n)=\wt(n)\bmod2\in\{0,1\},
  \qquad
  \eps_n=(-1)^{\wt(n)}=1-2\tau(n)\in\{1,-1\}.              \label{eq:normalizations}
\end{equation}
Thus
\begin{equation}
  \tau(n)=\frac{1-\eps_n}{2},
  \qquad
  \eps_n=1-2\tau(n).                                       \label{eq:normalization-bridge}
\end{equation}
The signed normalization is the natural one for products, finite differences, and Fourier
analysis; the bit normalization is the natural one for substitutions and digit expansions.
Keeping both visible prevents a large class of otherwise easy sign mistakes.

For bitwise exclusive-or we write \(a\xor b\), and \(k\submask n\) means that every
one-bit of \(k\) is also a one-bit of \(n\).  The symbol \(\res{x}{m}\) denotes the least
nonnegative residue of an integer \(x\) modulo \(m\).  Empty sums are zero and empty
products are one.

\subsection{The original Pascal-row logarithm}

The formula that motivated the drafts hides the binary expansion inside one row of Pascal's
triangle.  Put
\begin{equation}
  \mathcal X(n)=\sum_{k=0}^{n}(-1)^{\binom nk}.              \label{eq:source-X}
\end{equation}
If \(\nu(n)\) denotes the number of odd entries in row \(n\), then each odd entry
contributes \(-1\) and each even entry contributes \(+1\), so
\begin{equation}
  \mathcal X(n)=n+1-2\nu(n).                                \label{eq:X-via-odd-count}
\end{equation}
Lucas's theorem gives \(\nu(n)=2^{\wt(n)}\), and hence
\begin{equation}
  n+1-\mathcal X(n)=2^{\wt(n)+1}.                            \label{eq:source-exact-power}
\end{equation}
The logarithm in the resulting formula is therefore exact, not an approximation:
\begin{equation}
  \boxed{
  \tau(n)=\frac{1+(-1)^{\log_2\!\left(n+1-
       \sum_{k=0}^{n}(-1)^{\binom nk}\right)}}{2}.}          \label{eq:source-binomial-log}
\end{equation}
The purpose of the atlas is not to restate \eqref{eq:source-binomial-log} many times, but to
show that it is one manifestation of a single parity character appearing in several
apparently unrelated structures.

\begin{boxedremark}
\textbf{Proof-status boundary.}
The repository already formalizes the binary recursion, the finite block product, Prouhet
cancellation and its sharp first surviving moment, iterated prefix sums, and the
binomial--logarithm formula.  The additional valuation, Pascal-residue, determinant,
rarefied-sum, spectral, and analytic formulae are proved or cited here as ordinary
mathematics; they are not silently promoted to existing Lean theorems.
\end{boxedremark}

\subsection{A compact formula index}

The following table previews representative identities.  Their assumptions, proofs, and
interrelations are developed later.

\renewcommand{\arraystretch}{1.18}
\begin{longtable}{@{}p{0.255\textwidth}p{0.67\textwidth}@{}}
\toprule
Viewpoint & Representative identity \\
\midrule
\endfirsthead
\toprule
Viewpoint & Representative identity \\
\midrule
\endhead
Binary recursion & \(\eps_{2n}=\eps_n\), \(\eps_{2n+1}=-\eps_n\). \\
Block concatenation & \(\eps_{2^m q+r}=\eps_q\eps_r\) for \(0\le r<2^m\). \\
Boolean character & \(\eps_{a\xor b}=\eps_a\eps_b\). \\
Digit product & \(\eps_n=\prod_{j\ge0}(1-2b_j(n))\). \\
Floor sum & \(\eps_n=(-1)^{n-\sum_{j\ge1}\lfloor n/2^j\rfloor}\). \\
Factorial valuation & \(\eps_n=(-1)^{n-\vtwo(n!)}\). \\
Central binomial & \(\eps_n=(-1)^{\vtwo\binom{2n}{n}}\). \\
Carry cocycle & \(\eps_{a+b}=\eps_a\eps_b(-1)^{\vtwo\binom{a+b}{a}}\). \\
Sparse Pascal columns & \(\eps_n=(-1)^{\sum_j\binom{n}{2^j}}\). \\
Pascal support & \(\nu(n)=2^{\wt(n)}\), hence \(\eps_n=(-1)^{\log_2\nu(n)}\). \\
Log-free Pascal form & \(\tau(n)=\res{\nu(n)}3-1\). \\
Finite product & \(\sum_{n<2^m}\eps_nz^n=\prod_{j<m}(1-z^{2^j})\). \\
Mahler equation & \(E(z)=(1-z)E(z^2)\). \\
Ruler convolution & \(n\eps_n=-\sum_{k=1}^n(2^{\vtwo(k)+1}-1)\eps_{n-k}\). \\
Prouhet cancellation & \(\sum_{n<2^m}\eps_np(n)=0\) for \(\deg p<m\). \\
Sharp moment & \(\sum_{n<2^m}\eps_nn^m=(-1)^m m!2^{\binom m2}\). \\
Prefix discrepancy & \(\sum_{n<N}\eps_n=0\) for even \(N\), and \(\eps_{(N-1)/2}\) for odd \(N\). \\
Walsh spectrum & The level-\(m\) block is one Boolean frequency. \\
Riesz product & \(2^{-m}|P_m(\e^{2\pi\ii x})|^2=\prod_{j<m}(1-\cos(2\pi2^jx))\). \\
Dirichlet transform & \(\sum_{n\ge0}\eps_n/(n+1)^s\) has an entire continuation. \\
Rarefied sum & Residue-class sums at \(2^m\) are explicit powers of \(3\). \\
\bottomrule
\end{longtable}
\renewcommand{\arraystretch}{1}

\section{Binary recursion, blocks, automata, and uniqueness}
\label{sec:binary}

\subsection{The two-scale recurrence}

\begin{theorem}[Binary recursion and uniqueness]
\label{thm:binary-recursion}
The signed and zero--one sequences satisfy
\begin{align}
  \eps_0&=1,& \eps_{2n}&=\eps_n,& \eps_{2n+1}&=-\eps_n,     \label{eq:eps-recursion}\\
  \tau(0)&=0,& \tau(2n)&=\tau(n),& \tau(2n+1)&=1-\tau(n).  \label{eq:tau-recursion}
\end{align}
Either row, with its initial value, uniquely determines its sequence.
\end{theorem}

\begin{proof}
Multiplication by two appends a zero binary digit, while passing from \(2n\) to \(2n+1\)
replaces that terminal zero by one.  Thus
\(\wt(2n)=\wt(n)\) and \(\wt(2n+1)=\wt(n)+1\), proving the recurrences.
Every positive integer is \(2q\) or \(2q+1\) with \(q<n\), so strong induction reduces any
value to the initial value at zero.
\end{proof}

The infinite zero--one word
\[
  \mathbf t=0110100110010110\cdots
\]
is the fixed point beginning in zero of the length-two morphism
\begin{equation}
  \mu(0)=01,
  \qquad
  \mu(1)=10.                                                \label{eq:TM-morphism}
\end{equation}
In signed form, every block is followed by its negative.  If
\(\bm\eps^{(m)}=(\eps_0,\ldots,\eps_{2^m-1})\), then
\begin{equation}
  \bm\eps^{(m+1)}=(\bm\eps^{(m)},-\bm\eps^{(m)}),
  \qquad
  \bm\eps^{(m)}=(1,-1)^{\otimes m}.                         \label{eq:tensor-block}
\end{equation}

\subsection{Concatenating more than one bit}

\begin{theorem}[Dyadic block multiplicativity]
\label{thm:block-multiplicativity}
For \(m,q\ge0\) and \(0\le r<2^m\),
\begin{align}
  \wt(2^m q+r)&=\wt(q)+\wt(r),                              \label{eq:weight-block}\\
  \eps_{2^m q+r}&=\eps_q\eps_r,                             \label{eq:eps-block}\\
  \tau(2^m q+r)&=\tau(q)\xor\tau(r).                       \label{eq:tau-block}
\end{align}
\end{theorem}

\begin{proof}
The binary digits of \(q\) are shifted above the lowest \(m\) positions, while \(r\) occupies
only those positions.  Their supports are disjoint, so the digit sums add without carries.
Taking parity gives the last two identities.
\end{proof}

For each fixed \(m,r\), the subsequence \((\eps_{2^m q+r})_{q\ge0}\) is therefore either the
original sequence or its negative.  Consequently the signed \(2\)-kernel has only two
members; in the bit normalization the kernel consists of \(\tau\) and \(1-\tau\).  This is
the two-state automaton content of \eqref{eq:eps-recursion}.

\subsection{Complement and permutation symmetries}

\begin{proposition}[Reflection in a dyadic block]
\label{prop:block-reflection}
For \(0\le r<2^m\),
\begin{equation}
  \eps_{2^m-1-r}=(-1)^m\eps_r,
  \qquad
  \tau(2^m-1-r)=\tau(r)\xor(m\bmod2).                       \label{eq:block-reflection}
\end{equation}
More generally, permuting the first \(m\) binary positions leaves \(\eps_r\) unchanged.
\end{proposition}

\begin{proof}
Within exactly \(m\) positions, \(2^m-1-r\) is the bitwise complement of \(r\), so its
weight is \(m-\wt(r)\).  A permutation of positions does not change their sum.
\end{proof}

Thus a dyadic word is a palindrome when \(m\) is even and an anti-palindrome when \(m\) is
odd.  This elementary identity later explains reciprocity of the finite generating
polynomial and reversal closure of the factor language.

\subsection{The Boolean group law}

\begin{theorem}[XOR character]
\label{thm:xor-character}
For all \(a,b\ge0\),
\begin{equation}
  \eps_{a\xor b}=\eps_a\eps_b,
  \qquad
  \tau(a\xor b)=\tau(a)\xor\tau(b).                         \label{eq:xor-character}
\end{equation}
\end{theorem}

\begin{proof}
At every bit, exclusive-or is addition in \(\Ftwo\).  Hence
\(\wt(a\xor b)\equiv\wt(a)+\wt(b)\pmod2\), and both identities follow.
\end{proof}

This is the cleanest structural description of the signed sequence: \(\eps\) is a
one-dimensional character of the countable Boolean group of finite binary strings.
Ordinary addition is more complicated only because it introduces carries; \Cref{sec:valuations}
will identify the exact correction factor.

\section{Digit, floor, fractional-part, and trigonometric formulae}
\label{sec:digits}

\subsection{Extracting digits arithmetically}

\begin{lemma}[Floor extraction]
For every \(j\ge0\),
\begin{equation}
  b_j(n)=\left\lfloor\frac{n}{2^j}\right\rfloor
        -2\left\lfloor\frac{n}{2^{j+1}}\right\rfloor.       \label{eq:floor-bit}
\end{equation}
Consequently
\begin{equation}
  \eps_n=\prod_{j\ge0}\left(1-2\left(
    \left\lfloor\frac{n}{2^j}\right\rfloor
    -2\left\lfloor\frac{n}{2^{j+1}}\right\rfloor\right)\right),
  \label{eq:floor-digit-product}
\end{equation}
where all but finitely many factors are one.
\end{lemma}

\begin{proof}
The integer \(\lfloor n/2^j\rfloor\) has least significant bit \(b_j(n)\); subtracting twice
its quotient by two extracts that bit.  Since \((-1)^b=1-2b\) for \(b\in\{0,1\}\), the
product is \((-1)^{\sum_jb_j(n)}\).
\end{proof}

Summing \eqref{eq:floor-bit} with weights is not the only floor representation.  The finite
identity
\begin{equation}
  \wt(n)=n-\sum_{j\ge1}\left\lfloor\frac{n}{2^j}\right\rfloor                  \label{eq:digit-sum-floor}
\end{equation}
is the base-two instance of Legendre's digit-sum formula.  It gives
\begin{equation}
  \boxed{
  \eps_n=(-1)^{\,n-\sum_{j\ge1}\lfloor n/2^j\rfloor}.}     \label{eq:eps-floor-sum}
\end{equation}

\subsection{Fractional parts}

Because \(\sum_{j\ge1}n/2^j=n\), subtracting the floor sum in
\eqref{eq:digit-sum-floor} gives an exact infinite formula:
\begin{equation}
  \boxed{
  \wt(n)=\sum_{j\ge1}\left\{\frac{n}{2^j}\right\}.}        \label{eq:fractional-weight}
\end{equation}
The series is not effectively finite, but its tail is geometric once \(2^j>n\).  Thus
\begin{equation}
  \eps_n=\exp\!\left(\pi\ii\sum_{j\ge1}\left\{\frac{n}{2^j}\right\}\right).
  \label{eq:fractional-exponential}
\end{equation}
Although the exponent in \eqref{eq:fractional-exponential} is presented as an infinite sum,
\eqref{eq:fractional-weight} proves that it is the integer \(\wt(n)\).

\subsection{Finite sign products}

For a nonintegral real \(x\), \(\sgn(\sin\pi x)=(-1)^{\lfloor x\rfloor}\).  Since
\(\lfloor(n+1/2)/2^j\rfloor=\lfloor n/2^j\rfloor\), summing the floor exponents yields
\(2n-\wt(n)\).  Therefore
\begin{equation}
  \boxed{
  \eps_n=\prod_{j\ge0}
  \sgn\!\sin\!\left(\frac{\pi(n+1/2)}{2^j}\right).}        \label{eq:sine-sign-product}
\end{equation}
Once \(2^j>n+1/2\), the sine argument lies in \((0,\pi)\), so the factor is \(+1\).
An equivalent cosine form is
\begin{equation}
  \boxed{
  \eps_n=\prod_{j\ge0}\cos\!\left(\pi\left\lfloor\frac n{2^j}\right\rfloor\right).}
  \label{eq:floor-cosine-product}
\end{equation}
These products are exact square-wave encodings; they do not smooth the sequence.  The
half-unit shift in \eqref{eq:sine-sign-product} is essential because it prevents zero sine
factors at integer arguments.

\section{Two-adic valuations, central binomial coefficients, Catalan numbers, and carries}
\label{sec:valuations}

\subsection{Legendre's formula as a parity extractor}

For a positive integer \(M\), let \(\vtwo(M)\) be the exponent of two in its prime
factorization.  Legendre's formula gives
\begin{equation}
  \vtwo(n!)=\sum_{j\ge1}\left\lfloor\frac n{2^j}\right\rfloor
           =n-\wt(n).                                      \label{eq:legendre-two}
\end{equation}
Hence
\begin{equation}
  \boxed{
  \eps_n=(-1)^{n-\vtwo(n!)},
  \qquad
  \tau(n)=\res{n-\vtwo(n!)}2.}                             \label{eq:factorial-formulas}
\end{equation}

Applying \eqref{eq:legendre-two} to \((2n)!/(n!)^2\) collapses the whole digit sum into one
central binomial coefficient:
\begin{equation}
  \vtwo\binom{2n}{n}=2\wt(n)-\wt(2n)=\wt(n).                \label{eq:central-valuation}
\end{equation}
Thus
\begin{equation}
  \boxed{
  \eps_n=(-1)^{\vtwo\binom{2n}{n}},
  \qquad
  \tau(n)=\sin^2\!\left(\frac\pi2\vtwo\binom{2n}{n}\right).}                  \label{eq:central-binomial-formula}
\end{equation}
This is arguably the most economical arithmetic replacement for the row-wide
binomial--logarithm formula: one valuation of one integer contains exactly \(\wt(n)\).

\subsection{Catalan parity and the ruler sequence}

Let \(C_n=(n+1)^{-1}\binom{2n}{n}\) be the \(n\)-th Catalan number.  From
\eqref{eq:central-valuation},
\begin{equation}
  \vtwo(C_n)=\wt(n)-\vtwo(n+1)=\wt(n+1)-1.                  \label{eq:catalan-valuation}
\end{equation}
Consequently
\begin{equation}
  \boxed{
  \eps_n=(-1)^{\vtwo(C_n)+\vtwo(n+1)},
  \qquad
  \eps_{n+1}=-(-1)^{\vtwo(C_n)}.}                          \label{eq:catalan-signs}
\end{equation}
The second formula says that the parity of a Catalan valuation is the shifted complement of
Thue--Morse.

Incrementing \(n\) flips its trailing ones to zero and the next zero to one.  The number of
trailing ones is \(\vtwo(n+1)\), so
\begin{equation}
  \wt(n+1)=\wt(n)+1-\vtwo(n+1).                             \label{eq:weight-increment}
\end{equation}
Exponentiating gives a sequential recurrence with no explicit digit extraction:
\begin{equation}
  \boxed{
  \eps_{n+1}=-(-1)^{\vtwo(n+1)}\eps_n.}                    \label{eq:ruler-transition}
\end{equation}
Thus the parity of the classical ruler function is the multiplicative first difference of
Thue--Morse.

\subsection{Kummer's carry cocycle}

Kummer's theorem, in its base-two digit-sum form, states
\begin{equation}
  \vtwo\binom{a+b}{a}=\wt(a)+\wt(b)-\wt(a+b).               \label{eq:kummer-digit}
\end{equation}
The valuation is the number of binary carries, counted with propagation.  Rearranging parity
in \eqref{eq:kummer-digit} gives the exact obstruction to multiplicativity under ordinary
addition.

\begin{theorem}[Carry-corrected addition]
\label{thm:carry-corrected}
For all \(a,b\ge0\),
\begin{equation}
  \boxed{
  \eps_{a+b}=\eps_a\eps_b(-1)^{\vtwo\binom{a+b}{a}}.}       \label{eq:carry-corrected}
\end{equation}
Equivalently,
\begin{equation}
  \tau(a+b)=\tau(a)\xor\tau(b)\xor
    \res{\vtwo\binom{a+b}{a}}2.                            \label{eq:carry-corrected-bit}
\end{equation}
\end{theorem}

\begin{proof}
Equation \eqref{eq:kummer-digit} rewrites \(\wt(a+b)\) as
\(\wt(a)+\wt(b)-\vtwo\binom{a+b}{a}\).  Raise \((-1)\) to both sides; subtraction and
addition coincide modulo two.
\end{proof}

In particular, if \(a\bitand b=0\), then addition is carry-free and
\begin{equation}
  \eps_{a+b}=\eps_a\eps_b.                                  \label{eq:carry-free}
\end{equation}
The dyadic block formula is the most important special case.  More generally, if
\(N=a_1+\cdots+a_r\), then
\begin{equation}
  \eps_N=\left(\prod_{j=1}^r\eps_{a_j}\right)
  (-1)^{\vtwo\binom{N}{a_1,\ldots,a_r}},                    \label{eq:multinomial-carry}
\end{equation}
proved by applying Legendre's formula to the multinomial coefficient.

\section{Pascal support, sparse columns, residue formulae, and Boolean inversion}
\label{sec:pascal}

\subsection{Lucas support is a Boolean lattice}

Lucas's theorem modulo two says
\begin{equation}
  \binom nk\equiv\prod_{j\ge0}\binom{b_j(n)}{b_j(k)}\pmod2. \label{eq:lucas-binary}
\end{equation}
Every factor is one exactly when \(b_j(k)\le b_j(n)\).  Hence
\begin{equation}
  \binom nk\text{ is odd}\quad\Longleftrightarrow\quad k\submask n.             \label{eq:odd-submask}
\end{equation}
The odd positions in row \(n\) are therefore literally the Boolean lattice on the one-bits
of \(n\).  Its support polynomial factors as
\begin{equation}
  Q_n(z)=\sum_{k=0}^{n}\ind_{\binom nk\text{ odd}}z^k
        =\prod_{j:b_j(n)=1}(1+z^{2^j}).                     \label{eq:row-support-polynomial}
\end{equation}
Setting \(z=1\) proves Glaisher's count
\begin{equation}
  \boxed{\nu(n)=2^{\wt(n)}.}                                \label{eq:glaisher-count}
\end{equation}

\subsection{Power-of-two columns recover the digits}

Taking \(k=2^j\) in Lucas's theorem gives
\begin{equation}
  \binom{n}{2^j}\equiv b_j(n)\pmod2.                        \label{eq:power-column-digit}
\end{equation}
Thus only logarithmically many Pascal entries are needed to reconstruct the parity bit:
\begin{equation}
  \boxed{
  \eps_n=(-1)^{\sum_{j\ge0}\binom{n}{2^j}}
  =\prod_{j\ge0}\cos\!\left(\pi\binom{n}{2^j}\right).}    \label{eq:sparse-pascal-formula}
\end{equation}
Terms with \(2^j>n\) are zero.  Formula \eqref{eq:sparse-pascal-formula} is a direct sparse
companion to \eqref{eq:source-binomial-log}: it reads the row at power-of-two columns instead
of counting all of its odd entries.

\subsection{Removing the logarithm modulo three}

Because \(2^r\) alternates between \(1\) and \(2\) modulo three, Glaisher's count itself
reveals the parity of \(\wt(n)\) without a logarithm.

\begin{theorem}[Log-free odd-count formula]
\label{thm:odd-count-mod3}
For every \(n\ge0\),
\begin{equation}
  \boxed{
  \tau(n)=\res{\nu(n)}3-1,
  \qquad
  \eps_n=3-2\res{\nu(n)}3.}                                \label{eq:odd-count-mod3}
\end{equation}
The residue \(\res{\nu(n)}3\) is always \(1\) or \(2\).
\end{theorem}

\begin{proof}
By \eqref{eq:glaisher-count}, \(\nu(n)=2^{\wt(n)}\).  Its residue is one when the exponent is
even and two when it is odd.
\end{proof}

A version using exactly the signs from the original formula is also available.  Define
\begin{equation}
  S(n)=\sum_{k=1}^{n}(-1)^{\binom nk},\qquad S(0)=0.         \label{eq:S-pascal}
\end{equation}
Since the omitted \(k=0\) term equals \(-1\), equations
\eqref{eq:X-via-odd-count} and \eqref{eq:glaisher-count} give
\begin{equation}
  S(n)=n+2-2^{\wt(n)+1}.                                    \label{eq:S-exact}
\end{equation}
Therefore
\begin{align}
  \boxed{\tau(n)=\res{2n+S(n)}3,}                           \label{eq:pascal-residue}\\
  \boxed{\tau(n)=\frac43\sin^2\!\left(\frac\pi3(n-S(n))\right).}             \label{eq:pascal-sine}
\end{align}
Indeed \(2n+S(n)\) has residue zero or one according as \(\wt(n)\) is even or odd, while
\(n-S(n)\) has residue zero or two.

\subsection{A signed row-statistic family}

For a parameter \(u\), put
\begin{equation}
  R_u(n)=\sum_{k=0}^{n}u^{\binom nk\bmod2}.                 \label{eq:Ru-def}
\end{equation}
Even entries contribute one and odd entries contribute \(u\), so
\begin{equation}
  \boxed{R_u(n)=n+1+(u-1)2^{\wt(n)}.}                       \label{eq:Ru-exact}
\end{equation}
The original statistic is \(R_{-1}(n)=\mathcal X(n)\); \(R_0(n)\) counts even entries; and
any fixed \(u\ne1\) recovers the odd count from
\begin{equation}
  2^{\wt(n)}=\frac{R_u(n)-(n+1)}{u-1}.                      \label{eq:Ru-recovery}
\end{equation}
This makes clear that the special-looking sign sum is one point of a simple affine family.

\subsection{An \texorpdfstring{\(r\)}{r}-nomial logarithm family}

Expand \((x_1+\cdots+x_r)^n\).  A multinomial coefficient
\(\binom{n}{k_1,\ldots,k_r}\) is odd precisely when the binary additions
\(k_1+\cdots+k_r=n\) are carry-free.  At each one-bit of \(n\), exactly one of the \(r\)
indices receives that bit, independently of all other one-bits.  Hence:

\begin{theorem}[Odd multinomial count]
\label{thm:odd-multinomial}
The number \(\nu_r(n)\) of odd coefficients in the full \(r\)-nomial expansion of degree
\(n\) is
\begin{equation}
  \boxed{\nu_r(n)=r^{\wt(n)}.}                              \label{eq:odd-multinomial-count}
\end{equation}
For every integer \(r\ge2\),
\begin{equation}
  \eps_n=(-1)^{\log_r\nu_r(n)},
  \qquad
  \tau(n)=\res{\log_r\nu_r(n)}2.                           \label{eq:rnomial-log}
\end{equation}
\end{theorem}

The binomial row is the case \(r=2\).  The family shows that the logarithm is not tied to
Pascal's triangle: it measures the number of independent choices attached to the one-bits.

\subsection{Boolean M\"obius inversion}

The interval \([k,n]\) in the submask poset is a Boolean lattice of rank
\(\wt(n)-\wt(k)\).  Its M\"obius function is
\begin{equation}
  \mu(k,n)=(-1)^{\wt(n)-\wt(k)}=\eps_n\eps_k
  \qquad(k\submask n).                                     \label{eq:boolean-mobius}
\end{equation}
Therefore:

\begin{theorem}[Thue--Morse M\"obius inversion]
\label{thm:boolean-mobius}
For functions \(F,G:\N\to A\) valued in an abelian group, the relations
\begin{equation}
  F(n)=\sum_{k\submask n}G(k)                               \label{eq:boolean-zeta}
\end{equation}
are equivalent to
\begin{equation}
  \boxed{G(n)=\eps_n\sum_{k\submask n}\eps_kF(k).}         \label{eq:boolean-inversion}
\end{equation}
\end{theorem}

In matrix form, let
\(P_{n,k}=\ind_{\binom nk\text{ odd}}\) for \(0\le k\le n\).  Every finite lower-triangular
truncation has inverse
\begin{equation}
  (P^{-1})_{n,k}=\eps_n\eps_kP_{n,k}.                        \label{eq:pascal-parity-inverse}
\end{equation}
Thus the Thue--Morse signs are exactly the M\"obius kernel of Pascal's triangle reduced to
its parity support.

\section{Odious and evil enumerators}
\label{sec:odious}

An integer is called \emph{evil} when its binary weight is even and \emph{odious} when its
binary weight is odd.  Let \(\mathcal E_n\) and \(\mathcal O_n\) be the \(n\)-th evil and
odious nonnegative integers, respectively, with both sequences indexed from zero.

\begin{theorem}[Exact evil and odious enumerators]
\label{thm:evil-odious-enumerators}
For every \(n\ge0\),
\begin{equation}
  \boxed{
  \mathcal E_n=2n+\tau(n),
  \qquad
  \mathcal O_n=2n+1-\tau(n).}                               \label{eq:evil-odious-enumerators}
\end{equation}
Consequently
\begin{equation}
  \mathcal O_n-\mathcal E_n=\eps_n,
  \qquad
  \mathcal E_n+\mathcal O_n=4n+1.                          \label{eq:evil-odious-pair}
\end{equation}
\end{theorem}

\begin{proof}
The pair \(2n,2n+1\) has weights \(\wt(n)\) and \(\wt(n)+1\), so it contains exactly one
evil and one odious number.  If \(\tau(n)=0\), the even member \(2n\) is evil; if
\(\tau(n)=1\), the odd member \(2n+1\) is evil.  This gives both enumerators, and
\eqref{eq:evil-odious-pair} follows from \(\eps_n=1-2\tau(n)\).
\end{proof}

The formulas turn the sequence into a displacement between two complementary enumerations.
They also give immediate counting functions.  Let
\begin{equation}
  D(N)=\sum_{n=0}^{N-1}\eps_n.                              \label{eq:prefix-D-preview}
\end{equation}
Then the numbers of evil and odious integers below \(N\) are
\begin{equation}
  E(N)=\frac{N+D(N)}2,
  \qquad
  O(N)=\frac{N-D(N)}2.                                      \label{eq:evil-odious-counts}
\end{equation}
The exact prefix formula in \Cref{sec:prefix} will show \(|D(N)|\le1\).  Thus every even
initial interval contains exactly as many evil as odious integers, and every odd initial
interval differs by exactly one.

\section{Lacunary generating products and their coefficient formulae}
\label{sec:generating}

\subsection{The finite mother product}

For \(m\ge0\), define the level-\(m\) block polynomial
\begin{equation}
  P_m(z)=\sum_{n=0}^{2^m-1}\eps_nz^n.                       \label{eq:Pm-definition}
\end{equation}

\begin{theorem}[Finite lacunary product]
\label{thm:finite-lacunary-product}
For every \(m\ge0\),
\begin{equation}
  \boxed{
  P_m(z)=\prod_{j=0}^{m-1}(1-z^{2^j}).}                    \label{eq:finite-lacunary-product}
\end{equation}
In particular,
\begin{equation}
  P_{m+1}(z)=(1-z^{2^m})P_m(z).                             \label{eq:Pm-recursion}
\end{equation}
\end{theorem}

\begin{proof}
Expanding the product amounts to choosing a subset \(S\subseteq\{0,\ldots,m-1\}\).  It
contributes
\((-1)^{|S|}z^{\sum_{j\in S}2^j}\).  Binary representation gives a unique subset for each
\(0\le n<2^m\), and the coefficient is \((-1)^{\wt(n)}=\eps_n\).
\end{proof}

The product is the common source of the finite-difference, Prouhet, Fourier, and Riesz
formulae below.  It also displays exact reciprocity and cyclotomic multiplicities.

\begin{proposition}[Reciprocity and cyclotomic structure]
\label{prop:Pm-cyclotomic}
Let \(D=2^m-1\).  Then
\begin{equation}
  z^DP_m(z^{-1})=(-1)^mP_m(z),                              \label{eq:Pm-reciprocity}
\end{equation}
and
\begin{equation}
  \boxed{
  P_m(z)=(-1)^m\prod_{q=0}^{m-1}\Phi_{2^q}(z)^{m-q},}      \label{eq:Pm-cyclotomic}
\end{equation}
where \(\Phi_1(z)=z-1\).  A root of unity of exact order \(2^q\), with
\(0\le q<m\), has multiplicity \(m-q\).
\end{proposition}

\begin{proof}
For reciprocity,
\[
  z^DP_m(z^{-1})=\prod_{j<m}(z^{2^j}-1)=(-1)^mP_m(z).
\]
Also \(z^{2^j}-1=\prod_{q=0}^{j}\Phi_{2^q}(z)\).  Since
\(1-z^{2^j}=-(z^{2^j}-1)\), multiplying over \(j<m\) gives the sign and the exponent
\(m-q\) of each cyclotomic factor.
\end{proof}

Coefficient comparison in \eqref{eq:Pm-reciprocity} recovers the reflection formula
\eqref{eq:block-reflection}.  The zero of order \(m\) at \(z=1\) is the generating-function
shadow of Prouhet cancellation.

\subsection{The Mahler product}

The stabilized series
\begin{equation}
  E(z)=\sum_{n\ge0}\eps_nz^n                              \label{eq:E-definition}
\end{equation}
exists coefficientwise as a formal series.  For \(|z|<1\), the product converges normally
on compact subsets because \(\sum_j|z|^{2^j}<\infty\).  Passing to the limit in
\eqref{eq:finite-lacunary-product} gives
\begin{equation}
  \boxed{
  E(z)=\prod_{j\ge0}(1-z^{2^j}),
  \qquad
  E(z)=(1-z)E(z^2).}                                       \label{eq:E-Mahler}
\end{equation}
The second identity is the first-order Mahler functional equation obtained by removing the
\(j=0\) factor and reindexing.

For the bit series
\begin{equation}
  T(z)=\sum_{n\ge0}\tau(n)z^n,                             \label{eq:T-definition}
\end{equation}
we have
\begin{align}
  T(z)&=\frac12\left(\frac1{1-z}-E(z)\right),              \label{eq:T-from-E}\\
  T(z)&=(1-z)T(z^2)+\frac{z}{1-z^2}.                       \label{eq:T-Mahler}
\end{align}
The first identity is the generating form of \(\eps_n=1-2\tau(n)\); the second can also be
proved directly by separating even and odd indices and using \eqref{eq:tau-recursion}.

\subsection{The logarithm reveals the ruler sequence}

Set
\begin{equation}
  a_k=\sum_{2^j\mid k}2^j=1+2+\cdots+2^{\vtwo(k)}
     =2^{\vtwo(k)+1}-1.                                    \label{eq:ak-ruler}
\end{equation}
For \(|z|<1\), absolute convergence permits the rearrangement
\begin{align}
  \log E(z)
  &=\sum_{j\ge0}\log(1-z^{2^j})
    =-\sum_{j\ge0}\sum_{r\ge1}\frac{z^{r2^j}}r            \notag\\
  &=-\sum_{k\ge1}\frac{a_k}{k}z^k.                        \label{eq:log-E}
\end{align}
Indeed, a term of degree \(k\) occurs for exactly the powers \(2^j\mid k\), and its
contribution is \(-2^j/k\).

Logarithmic differentiation gives an unusual all-previous-terms recurrence.

\begin{theorem}[Ruler-convolution recurrence]
\label{thm:ruler-convolution}
With \(a_k=2^{\vtwo(k)+1}-1\),
\begin{equation}
  \boxed{
  n\eps_n=-\sum_{k=1}^{n}a_k\eps_{n-k}\qquad(n\ge1),}      \label{eq:ruler-convolution}
\end{equation}
with \(\eps_0=1\).
\end{theorem}

\begin{proof}
Differentiate \eqref{eq:log-E}, multiply by
\(E(z)=\sum_{r\ge0}\eps_rz^r\), and compare the coefficient of \(z^{n-1}\).
\end{proof}

The weighted convolution on the right of \eqref{eq:ruler-convolution} is therefore always
exactly \(\pm n\).  Unlike the two-state recursion, this formula contains no digits of
\(n\); all binary information is hidden in the valuations of the ordinary indices
\(1,\ldots,n\).

\subsection{Partition and Bell-polynomial formulae}

Exponentiating \eqref{eq:log-E} and collecting terms of total degree \(n\) gives a closed
finite sum over integer partitions.

\begin{theorem}[Partition formula]
\label{thm:partition-formula}
For every \(n\ge0\),
\begin{equation}
  \boxed{
  \eps_n=\sum_{\substack{m_1,\ldots,m_n\ge0\\
                   m_1+2m_2+\cdots+nm_n=n}}
  \prod_{j=1}^{n}\frac1{m_j!}
       \left(-\frac{a_j}{j}\right)^{m_j},
  \qquad a_j=2^{\vtwo(j)+1}-1.}                            \label{eq:partition-formula}
\end{equation}
At \(n=0\), the unique empty choice contributes one.
\end{theorem}

\begin{proof}
Write
\[
  E(z)=\prod_{j\ge1}\exp\!\left(-\frac{a_j}{j}z^j\right)
      =\prod_{j\ge1}\sum_{m_j\ge0}
        \frac1{m_j!}\left(-\frac{a_j}{j}\right)^{m_j}z^{jm_j}.
\]
The coefficient of \(z^n\) receives exactly the choices with \(\sum_jjm_j=n\).
\end{proof}

Define the complete exponential Bell polynomials by
\begin{equation}
  \exp\!\left(\sum_{j\ge1}x_j\frac{z^j}{j!}\right)
   =\sum_{n\ge0}\Bell_n(x_1,\ldots,x_n)\frac{z^n}{n!}.     \label{eq:Bell-definition}
\end{equation}
Then \eqref{eq:log-E} yields the compact equivalent form
\begin{equation}
  \boxed{
  \eps_n=\frac1{n!}\Bell_n(-0!a_1,-1!a_2,\ldots,-(n-1)!a_n).}                  \label{eq:Bell-eps}
\end{equation}
The rational-looking partition sum and Bell polynomial always collapse to \(\pm1\); the
cancellation is highly nonlocal and invisible term by term.

\subsection{A lower-Hessenberg determinant}

The same recurrence admits a determinant solution.  For \(n\ge1\), let
\begin{equation}
 H_n=\begin{pmatrix}
 a_1&1&0&0&\cdots&0\\
 a_2&a_1&2&0&\cdots&0\\
 a_3&a_2&a_1&3&\ddots&\vdots\\
 \vdots&\vdots&\ddots&\ddots&\ddots&0\\
 a_{n-1}&a_{n-2}&\cdots&a_2&a_1&n-1\\
 a_n&a_{n-1}&\cdots&a_3&a_2&a_1
 \end{pmatrix}.                                            \label{eq:Hn-matrix}
\end{equation}

\begin{theorem}[Hessenberg determinant formula]
\label{thm:Hessenberg-determinant}
For every \(n\ge1\),
\begin{equation}
  \boxed{
  \eps_n=\frac{(-1)^n}{n!}\det H_n.}                       \label{eq:Hessenberg-determinant}
\end{equation}
In particular, \(|\det H_n|=n!\).
\end{theorem}

\begin{proof}
Let \(D_0=1\) and \(D_n=\det H_n\).  Expansion of a lower-Hessenberg determinant along the
last row, or induction by continuant elimination, gives
\begin{equation}
  D_n=\sum_{r=1}^{n}(-1)^{r+1}
      \frac{(n-1)!}{(n-r)!}\,a_rD_{n-r}.                    \label{eq:Hessenberg-recurrence}
\end{equation}
Put \(c_n=(-1)^nD_n/n!\).  Multiplying \eqref{eq:Hessenberg-recurrence} by
\((-1)^n/n!\) gives
\[
  nc_n=-\sum_{r=1}^{n}a_rc_{n-r},\qquad c_0=1.
\]
This is exactly \eqref{eq:ruler-convolution}, whose solution is unique by induction.
\end{proof}

The determinant is structurally different from the Hankel determinants discussed in
\Cref{sec:hankel}.  Here the matrix is engineered from the logarithmic derivative and its
determinant has the elementary magnitude \(n!\); the classical Hankel matrix uses the
Thue--Morse terms themselves and has a much subtler nonvanishing theorem.

\subsection{The bit series over characteristic two}

Reducing \eqref{eq:T-Mahler} coefficientwise modulo two, use
\(T(z^2)=T(z)^2\), \(1-z=1+z\), and \(1+z^2=(1+z)^2\).  One obtains
\begin{equation}
  \boxed{
  (1+z)^3T(z)^2+(1+z)^2T(z)+z=0
  \qquad\text{in }\Ftwo[[z]].}                              \label{eq:T-F2-algebraic}
\end{equation}
This is an explicit algebraic certificate of automaticity.  It is also a concrete instance
of the Christol--Kamae--Mend\`es France--Rauzy theorem: over a finite field, a sequence is
automatic exactly when its generating series is algebraic over the rational-function field
\cite{christol1980}.

\subsection{An integer algebraic lift}

A surprisingly different integer sequence has Thue--Morse parity.  Define \(c_0=0\),
\(c_1=1\), and for \(n\ge2\),
\begin{equation}
  nc_n=(5-2n)c_{n-1}+3(n-1)c_{n-2}+1.                       \label{eq:integer-lift-recurrence}
\end{equation}
Its first terms are
\[
  0,1,1,2,1,4,-2,13,-23,68,-164,439,-1146,3067,\ldots .
\]

\begin{theorem}[Algebraic integer lift]
\label{thm:integer-lift}
Let \(C(z)=\sum_{n\ge0}c_nz^n\).  Then
\begin{align}
  C(z)&=\frac{1}{2(1-z)}
    \left(\sqrt{\frac{1+3z}{1-z}}-1\right),                 \label{eq:C-lift-closed}\\
  (1-z)^3C(z)^2+(1-z)^2C(z)&=z.                             \label{eq:C-lift-algebraic}
\end{align}
All \(c_n\) are integers and
\begin{equation}
  \boxed{c_n\equiv\tau(n)\pmod2,
  \qquad \eps_n=(-1)^{c_n}.}                               \label{eq:C-lift-parity}
\end{equation}
\end{theorem}

\begin{proof}
Multiplying \eqref{eq:integer-lift-recurrence} by \(z^{n-1}\), summing, and using the
initial conditions gives a first-order differential equation whose solution is
\eqref{eq:C-lift-closed}; squaring gives \eqref{eq:C-lift-algebraic}.  Conversely, because
\(C(0)=0\), coefficient comparison in \eqref{eq:C-lift-algebraic} determines each new
coefficient from earlier integer coefficients, proving integrality and uniqueness.

Modulo two, \eqref{eq:C-lift-algebraic} becomes exactly
\eqref{eq:T-F2-algebraic}.  That equation has a unique solution in \(z\Ftwo[[z]]\): at each
degree the linear term contains the new coefficient with coefficient one, while the square
uses only earlier or lower-index data.  Hence \(C(z)\equiv T(z)\pmod2\).
\end{proof}

This lift is not a disguised binary weight: its coefficients grow, change sign, and are
nonautomatic over \(\Z\).  The parity survives because an algebraic equation over
\(\Z[[z]]\) reduces to the characteristic-two equation of the automatic sequence.

\subsection{A direct natural-boundary argument}

The unit circle is a natural boundary for \(E\).  One can invoke the P\'olya--Carlson theorem,
but the product gives a concrete proof.  Every root of unity whose order is a power of two
is a radial zero: if \(\zeta^{2^J}=1\), then for \(0<r<1\) the tail contains
\(\prod_{j\ge J}(1-r^{2^j})\), which tends to zero as \(r\uparrow1\).  Such roots are dense
on the unit circle.  If \(E\) admitted analytic continuation through any open boundary arc,
the continuation would vanish at a dense set in that arc and hence, by the identity theorem,
would vanish identically in a neighborhood.  This contradicts \(E(0)=1\).  Therefore no
boundary arc can be crossed.

\section{Finite differences, Prouhet cancellation, and exact moments}
\label{sec:prouhet}

\subsection{A master mixed-difference identity}

For \(h\in\R\), let \(T_hf(x)=f(x+h)\) and
\(\Delta_h=T_h-I\).  Translation operators commute, and the coefficient of
\(T_{nh}\) in \(\prod_{j<m}(I-T_{2^jh})\) is \(\eps_n\).  Hence:

\begin{theorem}[Mixed finite-difference identity]
\label{thm:mixed-difference}
For every function on which the displayed translations are defined,
\begin{equation}
  \boxed{
  \sum_{n=0}^{2^m-1}\eps_nf(x+nh)
   =\prod_{j=0}^{m-1}(I-T_{2^jh})f(x)
   =(-1)^m\Delta_h\Delta_{2h}\cdots\Delta_{2^{m-1}h}f(x).}  \label{eq:mixed-difference}
\end{equation}
\end{theorem}

For \(h>0\) and \(f\in C^m\), iterating
\(\Delta_af(x)=\int_0^af'(x+u)\dd u\) gives the integral form
\begin{equation}
  \boxed{
  \sum_{n<2^m}\eps_nf(x+nh)
   =(-1)^m\int_0^h\int_0^{2h}\cdots\int_0^{2^{m-1}h}
       f^{(m)}(x+u_0+\cdots+u_{m-1})
       \dd u_{m-1}\cdots\dd u_0.}                          \label{eq:mixed-difference-integral}
\end{equation}
This formula contains both exact cancellation and sign information.  For example, if
\(x>0\), \(h>0\), and \(s>0\), then
\begin{equation}
  \boxed{
  \sum_{n<2^m}\frac{\eps_n}{(x+nh)^s}
  =(s)_m\int_0^h\cdots\int_0^{2^{m-1}h}
  \frac{\dd u_{m-1}\cdots\dd u_0}
       {(x+u_0+\cdots+u_{m-1})^{s+m}}>0,}                  \label{eq:positive-reciprocal-block}
\end{equation}
where \((s)_m=s(s+1)\cdots(s+m-1)\).  The sign \((-1)^m\) in
\eqref{eq:mixed-difference-integral} cancels the sign of the \(m\)-th derivative of
\(y^{-s}\).

\subsection{Prouhet annihilation and sharp extraction}

If \(p\) is a polynomial of degree below \(m\), then its \(m\)-th derivative vanishes.
Equation \eqref{eq:mixed-difference-integral} therefore gives Prouhet's construction.

\begin{theorem}[Prouhet annihilation]
\label{thm:prouhet-annihilation}
For every polynomial \(p\) with \(\deg p<m\),
\begin{equation}
  \boxed{
  \sum_{n=0}^{2^m-1}\eps_np(n)=0.}                          \label{eq:prouhet-polynomial}
\end{equation}
Equivalently, for \(0\le r<m\),
\begin{equation}
  \sum_{n<2^m}\eps_nn^r=0.                                 \label{eq:prouhet-moments-zero}
\end{equation}
\end{theorem}

Since \(\eps_n=+1\) on the evil class and \(-1\) on the odious class,
\eqref{eq:prouhet-moments-zero} says that the two classes below \(2^m\) have equal sums of
like powers through degree \(m-1\).  At the next degree the cancellation is exact but no
longer zero.

\begin{theorem}[First surviving moment]
\label{thm:first-surviving-moment}
For every \(m\ge0\),
\begin{equation}
  \boxed{
  \sum_{n=0}^{2^m-1}\eps_nn^m=(-1)^m m!\,2^{\binom m2}.}    \label{eq:first-surviving-moment}
\end{equation}
More generally, if \(p(x)=a_mx^m+\cdots+a_0\), then
\begin{equation}
  \sum_{n<2^m}\eps_np(n)=a_m(-1)^m m!2^{\binom m2}.         \label{eq:leading-coefficient-extractor}
\end{equation}
\end{theorem}

\begin{proof}
Take \(h=1\) and \(f(x)=x^m\) in \eqref{eq:mixed-difference-integral}.  The derivative is
the constant \(m!\), and the rectangular integration volume is
\(1\cdot2\cdots2^{m-1}=2^{\binom m2}\).  Lower-degree terms vanish by
\Cref{thm:prouhet-annihilation}.
\end{proof}

This is stronger than a vanishing theorem: on the vector space of polynomials of degree at
most \(m\), a complete Thue--Morse block is an exact leading-coefficient functional.

\subsection{The finite exponential transform}

Define
\begin{equation}
  F_m(t)=\sum_{n=0}^{2^m-1}\eps_n\e^{nt}.                  \label{eq:Fm-exponential}
\end{equation}
Substituting \(z=\e^t\) in \eqref{eq:finite-lacunary-product} gives
\begin{equation}
  \boxed{
  F_m(t)=\prod_{j=0}^{m-1}(1-\e^{2^jt}).}                  \label{eq:Fm-product}
\end{equation}
Using \(1-\e^u=-2\e^{u/2}\sinh(u/2)\) and
\(\sum_{j<m}2^j=2^m-1\),
\begin{equation}
  F_m(t)=(-2)^m\e^{(2^m-1)t/2}
          \prod_{j=0}^{m-1}\sinh(2^{j-1}t).                \label{eq:Fm-sinh}
\end{equation}
After centering the block at \(c_m=(2^m-1)/2\),
\begin{equation}
  \boxed{
  \sum_{n<2^m}\eps_n\e^{(n-c_m)t}
   =(-2)^m\prod_{j=0}^{m-1}\sinh(2^{j-1}t).}               \label{eq:centered-exponential}
\end{equation}
The zero of exact order \(m\) at the origin simultaneously proves
\eqref{eq:prouhet-moments-zero} and \eqref{eq:first-surviving-moment}.

\subsection{All higher dyadic power sums}

Let
\begin{equation}
  M_{m,r}=\sum_{n=0}^{2^m-1}\eps_nn^r.                      \label{eq:Mmr-def}
\end{equation}
Since \(F_m(t)=\sum_{r\ge0}M_{m,r}t^r/r!\), expanding each factor in
\eqref{eq:Fm-product} gives a finite formula valid at every order.

\begin{theorem}[Composition formula for all moments]
\label{thm:all-moments-composition}
For \(r\ge0\),
\begin{equation}
  \boxed{
  M_{m,r}=(-1)^m r!
  \sum_{\substack{q_0,\ldots,q_{m-1}\ge1\\q_0+\cdots+q_{m-1}=r}}
  \frac{2^{\sum_{j=0}^{m-1}jq_j}}{q_0!\cdots q_{m-1}!}.}   \label{eq:all-moments-composition}
\end{equation}
The sum is empty, and hence zero, when \(r<m\).
\end{theorem}

\begin{proof}
In the \(j\)-th factor,
\[
  1-\e^{2^jt}=-\sum_{q_j\ge1}\frac{2^{jq_j}t^{q_j}}{q_j!}.
\]
Multiply the \(m\) series and compare the coefficient of \(t^r\).
\end{proof}

There is also a compact Bell-polynomial form.  Let \(B_{2r}\) be the Bernoulli numbers, with
\(B_2=1/6\), and define
\begin{align}
  x_1^{(m)}&=\frac{2^m-1}{2},                               \label{eq:moment-cumulant-1}\\
  x_{2r}^{(m)}&=\frac{B_{2r}}{2r}\,
       \frac{4^{rm}-1}{4^r-1}\qquad(r\ge1),                \label{eq:moment-cumulant-even}\\
  x_{2r+1}^{(m)}&=0\qquad(r\ge1).                          \label{eq:moment-cumulant-odd}
\end{align}
The expansion
\[
  \log\frac{\sinh y}{y}
   =\sum_{r\ge1}\frac{2^{2r-1}B_{2r}}{r(2r)!}y^{2r}
\]
turns \eqref{eq:Fm-sinh} into
\begin{equation}
  F_m(t)=(-1)^m2^{\binom m2}t^m
  \exp\!\left(\sum_{k\ge1}x_k^{(m)}\frac{t^k}{k!}\right). \label{eq:Fm-Bell-normalized}
\end{equation}
Therefore:

\begin{corollary}[Bell-polynomial moments]
\label{cor:Bell-moments}
For every \(\ell\ge0\),
\begin{equation}
  \boxed{
  M_{m,m+\ell}=(-1)^m2^{\binom m2}
   \frac{(m+\ell)!}{\ell!}
   \Bell_\ell\!\left(x_1^{(m)},\ldots,x_\ell^{(m)}\right).}                   \label{eq:Bell-moments}
\end{equation}
\end{corollary}

For example,
\begin{align}
  M_{m,m+1}
   &=(-1)^m2^{\binom m2}(m+1)!\frac{2^m-1}{2},              \label{eq:moment-mplus1}\\
  M_{m,m+2}
   &=(-1)^m2^{\binom m2}(m+2)!
     \left(\frac{(2^m-1)^2}{8}+\frac{4^m-1}{72}\right).    \label{eq:moment-mplus2}
\end{align}
Thus the familiar first surviving moment is only the first coefficient of a complete exact
hierarchy.

\subsection{Sparse Prouhet identities on one Pascal row}

The usual Prouhet block is the full submask set of \(2^m-1\).  Lucas's criterion suggests
replacing it by an arbitrary Pascal row.  Suppose the one-bit positions of \(n\) are
\(a_1<\cdots<a_r\).  Then
\begin{equation}
  \sum_{k\submask n}\eps_kz^k
   =\prod_{j=1}^{r}(1-z^{2^{a_j}}).                         \label{eq:sparse-Prouhet-product}
\end{equation}
The right side has a zero of exact order \(r=\wt(n)\) at \(z=1\).

\begin{theorem}[Sparse Pascal-row Prouhet theorem]
\label{thm:sparse-Prouhet}
Let \(n\) have one-bit positions \(a_1,\ldots,a_r\).  Then
\begin{align}
  \sum_{k\submask n}\eps_kk^d&=0 &&(0\le d<r),             \label{eq:sparse-Prouhet-zero}\\
  \sum_{k\submask n}\eps_kk^r
    &=(-1)^r r!\,2^{a_1+\cdots+a_r}.                       \label{eq:sparse-Prouhet-boundary}
\end{align}
Equivalently, the odd support of row \(n\) of Pascal's triangle, signed by Thue--Morse,
annihilates every polynomial of degree below the number of one-bits of \(n\).
\end{theorem}

\begin{proof}
Set \(z=\e^t\) in \eqref{eq:sparse-Prouhet-product}.  Near zero,
\[
  \prod_{j=1}^r(1-\e^{2^{a_j}t})
   =(-1)^r2^{a_1+\cdots+a_r}t^r+\bigO(t^{r+1}).
\]
Compare exponential generating coefficients.
\end{proof}

\begin{example}
For \(n=13=(1101)_2\), the one-bit positions are \(0,2,3\).  The odd positions in Pascal
row 13 are the submasks
\(0,1,4,5,8,9,12,13\).  Their Thue--Morse-signed sums of powers vanish through degree two,
and
\[
  \sum_{k\submask13}\eps_kk^3=(-1)^3\,3!\,2^{0+2+3}=-192.
\]
\end{example}

\section{Exact prefix sums, discrepancy, and the Fabius connection}
\label{sec:prefix}

\subsection{The first prefix is already exact}

Define the half-open signed prefix sum
\begin{equation}
  D(N)=\sum_{n=0}^{N-1}\eps_n.                              \label{eq:D-prefix}
\end{equation}
Pairing consecutive terms gives a complete formula.

\begin{theorem}[Exact prefix discrepancy]
\label{thm:exact-prefix}
For every \(N\ge0\),
\begin{equation}
  \boxed{
  D(N)=
  \begin{cases}
    0,&N\text{ even},\\
    \eps_{(N-1)/2},&N\text{ odd}.
  \end{cases}}                                             \label{eq:exact-prefix}
\end{equation}
In particular, \(|D(N)|\le1\).
\end{theorem}

\begin{proof}
For \(N=2M\), group the sum into pairs
\(\eps_{2q}+\eps_{2q+1}=\eps_q-\eps_q=0\).  For \(N=2M+1\), the same pairs leave
\(\eps_{2M}=\eps_M\).
\end{proof}

Thus every interval has uniformly bounded discrepancy:
\begin{equation}
  \left|\sum_{n=A}^{B-1}\eps_n\right|
  =|D(B)-D(A)|\le2.                                        \label{eq:interval-discrepancy}
\end{equation}
Every aligned dyadic block of positive length has sum zero:
\begin{equation}
  \sum_{r=0}^{2^m-1}\eps_{q2^m+r}
  =\eps_q\sum_{r<2^m}\eps_r=0\qquad(m\ge1).               \label{eq:aligned-block-zero}
\end{equation}
This is stronger than asymptotic equidistribution; it is exact local balance at every dyadic
scale.

\subsection{Iterated prefixes and Prouhet vanishing}

Let \((\mathcal Pa)(N)=\sum_{0\le n<N}a(n)\), and let \(\mathcal P^r\) denote \(r\)
iterations.  For \(r\ge1\), repeated summation gives
\begin{equation}
  (\mathcal P^ra)(N)
   =\sum_{n=0}^{N-1}\binom{N-n-1}{r-1}a(n).                 \label{eq:iterated-prefix-binomial}
\end{equation}
For fixed \(N=2^m\), the binomial weight is a polynomial in \(n\) of degree \(r-1\).
Consequently
\begin{equation}
  (\mathcal P^r\eps)(2^m)=0
  \qquad(1\le r\le m).                                    \label{eq:iterated-prefix-dyadic-zero}
\end{equation}
Different inclusive/exclusive endpoint conventions shift this formula by one and account for
several superficially different signs in the literature and in computational experiments.
The invariant content is the same: an \(m\)-bit block kills all prefix kernels of polynomial
degree below \(m\).

\subsection{Why the sequence appears in the Fabius theory}

Let \(\Up\) denote Rvachev's even compactly supported up-function.  In the main
Fabius--Rvachev development, the signed global extension is represented by the locally
finite translate sum
\begin{equation}
  \boxed{
  \mathcal F(x)=\sum_{n\ge0}\eps_n\Up(x-2n-1).}            \label{eq:Fabius-global-extension}
\end{equation}
Because \(\Up\) is supported on \([-1,1]\), only finitely many terms contribute at any fixed
\(x\).  The alternating binary block structure is exactly what converts the refinement law
for \(\Up\) into the pure global dilation equation
\begin{equation}
  \mathcal F'(x)=2\mathcal F(2x).                           \label{eq:Fabius-global-dilation}
\end{equation}

The same signs govern exact dyadic evaluations of the bounded Fabius distribution function.
Finite spline and moment formulae produce sums over complete or truncated dyadic blocks;
\Cref{thm:prouhet-annihilation} removes all polynomial terms below the critical degree, while
\Cref{thm:first-surviving-moment} supplies the first nonzero contribution with its exact
power of two and factorial.  The block product \eqref{eq:finite-lacunary-product} is the
coefficient form of that cancellation.  Thus the Thue--Morse layer is not decorative: it is
the arithmetic mechanism that makes the dyadic formulas collapse.

The relevant formal modules include \lean{ThueMorseGenerating.lean},
\lean{ThueMorseExponential.lean}, \lean{ThueMorsePrefix.lean},
\lean{ThueMorseApproximation.lean}, \lean{PaperKFoldThueMorse.lean}, and the exact dyadic
modules listed in \Cref{sec:formalization}.  Formula \eqref{eq:source-binomial-log} is
formalized separately in \lean{ThueMorseBinomialLog.lean}.

\section{Ordinary Fourier transforms, Walsh characters, and Riesz products}
\label{sec:fourier}

\subsection{The trigonometric factorization}

On the unit circle, \eqref{eq:finite-lacunary-product} becomes a finite sine product.  Since
\(1-\e^{\ii u}=-2\ii\e^{\ii u/2}\sin(u/2)\),
\begin{equation}
  \boxed{
  \sum_{n<2^m}\eps_n\e^{\ii n\theta}
  =(-2\ii)^m\e^{\ii(2^m-1)\theta/2}
     \prod_{j=0}^{m-1}\sin(2^{j-1}\theta).}                \label{eq:finite-sine-product}
\end{equation}
The phase records the center of the block; the magnitude factors scale by scale:
\begin{equation}
  \left|\sum_{n<2^m}\eps_n\e^{\ii n\theta}\right|
  =2^m\prod_{j=0}^{m-1}|\sin(2^{j-1}\theta)|.              \label{eq:finite-sine-magnitude}
\end{equation}
This is the oscillatory counterpart of the centered hyperbolic-sine formula
\eqref{eq:centered-exponential}.

\subsection{The exact discrete Fourier transform}

Let \(N=2^m\), and use the convention
\begin{equation}
  \widehat\eps_m(k)=\sum_{n=0}^{N-1}\eps_n\e^{-2\pi\ii kn/N}
  \qquad(0\le k<N).                                        \label{eq:DFT-definition}
\end{equation}
Evaluation of the block polynomial gives
\begin{equation}
  \boxed{
  \widehat\eps_m(k)=
  \prod_{j=0}^{m-1}\left(1-\e^{-2\pi\ii k2^j/N}\right).}  \label{eq:exact-DFT}
\end{equation}
If \(k\ne0\) is even, one factor vanishes: writing \(k=2^vu\) with \(u\) odd, take
\(j=m-v\).  Hence, for \(m\ge1\), the ordinary spectrum is supported on odd frequencies:
\begin{equation}
  \widehat\eps_m(k)=0\qquad(k\text{ even}).                \label{eq:DFT-even-zero}
\end{equation}
For odd \(k\), every factor is nonzero, and Fourier inversion gives a finite cyclotomic
formula for each individual sign:
\begin{equation}
  \boxed{
  \eps_n=\frac1{2^m}\sum_{\substack{0\le k<2^m\\k\text{ odd}}}
  \left[\prod_{j=0}^{m-1}
     \left(1-\e^{-2\pi\ii k2^j/2^m}\right)\right]
  \e^{2\pi\ii kn/2^m},\qquad 0\le n<2^m.}                 \label{eq:DFT-inversion}
\end{equation}
This is computationally inferior to the digit formula, but it is the natural exact
representation for cyclic spectral questions.

\subsection{The Walsh transform is a delta function}

Identify \(0\le n<2^m\) with its bit vector
\(\bm b(n)=(b_0(n),\ldots,b_{m-1}(n))\in\Ftwo^m\).  For
\(\bm a\in\Ftwo^m\), let
\begin{equation}
  \chi_{\bm a}(\bm b)=(-1)^{\bm a\cdot\bm b}.              \label{eq:Walsh-character}
\end{equation}
The Thue--Morse sign is the top-order character
\(\chi_{\bm1}\), where \(\bm1=(1,\ldots,1)\).

\begin{theorem}[One-point Walsh spectrum]
\label{thm:Walsh-delta}
For \(0\le r<2^m\),
\begin{equation}
  \boxed{
  \sum_{n=0}^{2^m-1}(-1)^{\wt(n)+\wt(r\bitand n)}
  =\begin{cases}
     2^m,&r=2^m-1,\\
     0,&r\ne2^m-1.
   \end{cases}}                                            \label{eq:Walsh-delta}
\end{equation}
\end{theorem}

\begin{proof}
The summand is the product \(\chi_{\bm1}(\bm b(n))\chi_{\bm b(r)}(\bm b(n))\), hence the
character \(\chi_{\bm1+\bm b(r)}\).  A nontrivial character sums to zero over \(\Ftwo^m\),
while the trivial character sums to \(2^m\).
\end{proof}

Thus Thue--Morse is spectrally trivial for XOR: it is one Walsh frequency.  The rich ordinary
Fourier behavior comes entirely from replacing Boolean addition by cyclic addition.

\subsection{Finite Riesz products}

Normalize the ordinary Fourier intensity by
\begin{equation}
  R_m(x)=2^{-m}\left|P_m(\e^{2\pi\ii x})\right|^2.          \label{eq:Riesz-def}
\end{equation}
Since \(|1-\e^{\ii u}|^2=2(1-\cos u)\),
\begin{equation}
  \boxed{
  R_m(x)=\prod_{j=0}^{m-1}\left(1-\cos(2\pi2^jx)\right).}  \label{eq:Riesz-product}
\end{equation}
Parseval's identity, or the constant coefficient of the product, gives
\begin{equation}
  \int_0^1R_m(x)\dd x=1.                                   \label{eq:Riesz-mass-one}
\end{equation}
Thus \(R_m(x)\dd x\) is a probability measure.  It does not converge to a smooth density;
its weak limit is the classical singular continuous Thue--Morse spectral measure
\cite{baakegrimm2008,queffelec2010}.  Formula \eqref{eq:Riesz-product} is the exact finite
object from which that measure is built.

\subsection{Autocorrelation recursion}

Define the limiting autocorrelation, whose existence follows from unique ergodicity of the
primitive substitution system, by
\begin{equation}
  \eta(r)=\lim_{N\to\infty}\frac1N
  \sum_{n=0}^{N-1}\eps_n\eps_{n+r}\qquad(r\ge0).           \label{eq:autocorrelation-def}
\end{equation}
Separating even and odd indices and using \eqref{eq:eps-recursion} gives
\begin{equation}
  \boxed{
  \eta(0)=1,
  \qquad
  \eta(2r)=\eta(r),
  \qquad
  \eta(2r+1)=-\frac{\eta(r)+\eta(r+1)}2.}                  \label{eq:autocorrelation-recursion}
\end{equation}
For example, the odd recurrence at \(r=0\) gives
\(\eta(1)=-\tfrac12(1+\eta(1))\), hence
\begin{equation}
  \eta(1)=-\frac13,
  \qquad \eta(2)=-\frac13,
  \qquad \eta(3)=\frac13.                                 \label{eq:autocorrelation-first}
\end{equation}
The numbers \(\eta(r)\) are the Fourier coefficients of the limiting Riesz-product measure.
The recursion is elementary; the theorem that the measure is purely singular continuous is
a deeper spectral result.

\section{Dirichlet transforms, digit constants, and infinite products}
\label{sec:analytic}

\subsection{An entire shifted Dirichlet series}

Consider
\begin{equation}
  \mathcal D(s)=\sum_{n\ge0}\frac{\eps_n}{(n+1)^s}.        \label{eq:Dirichlet-definition}
\end{equation}
The bounded prefix sums from \Cref{thm:exact-prefix} imply convergence for
\(\Re s>0\) by summation by parts, and absolute convergence for \(\Re s>1\).  In the latter
half-plane, the Mellin formula for \((n+1)^{-s}\) gives
\begin{equation}
  \Gamma(s)\mathcal D(s)
   =\int_0^\infty t^{s-1}\e^{-t}E(\e^{-t})\dd t.           \label{eq:Dirichlet-Mellin}
\end{equation}

\begin{theorem}[Entire continuation]
\label{thm:Dirichlet-entire}
The function \(\mathcal D(s)\) extends to an entire function on \(\C\).  At every
nonnegative integer \(r\),
\begin{equation}
  \boxed{\mathcal D(-r)=0.}                                \label{eq:Dirichlet-trivial-zeros}
\end{equation}
For real \(s>0\), \(\mathcal D(s)>0\).
\end{theorem}

\begin{proof}
As \(t\to\infty\), the factor \(\e^{-t}\) gives exponential decay.  As \(t\downarrow0\),
choose \(J\) with \(2^Jt\le1<2^{J+1}t\).  For \(j\le J\),
\(1-\e^{-2^jt}\le2^jt\), so
\[
  0<E(\e^{-t})
  \le t^{J+1}2^{J(J+1)/2}.
\]
Since \(J\asymp\log(1/t)/\log2\), this bound is smaller than every fixed power of \(t\).
Therefore the Mellin integral in \eqref{eq:Dirichlet-Mellin} defines an entire function of
\(s\).  Multiplication by the entire function \(1/\Gamma(s)\) gives the entire continuation
of \(\mathcal D\); the zeros of \(1/\Gamma\) at \(0,-1,-2,\ldots\) give
\eqref{eq:Dirichlet-trivial-zeros}.  Positivity follows because every factor of
\(E(\e^{-t})\) is positive for \(t>0\).
\end{proof}

The vanishing values are the analytic continuation of Prouhet cancellation: formally,
\(\mathcal D(-r)\) would be a divergent sum of \(\eps_n(n+1)^r\), and the continuation
assigns zero to every nonnegative polynomial degree.

\subsection{Thue--Morse digit constants in every base}

For an integer \(b\ge2\), define the base-\(b\) number with Thue--Morse digits
\begin{equation}
  \kappa_b=\sum_{n\ge0}\frac{\tau(n)}{b^{n+1}}.            \label{eq:kappa-b}
\end{equation}
Evaluating \eqref{eq:T-from-E} at \(z=1/b\) gives
\begin{equation}
  \boxed{
  \kappa_b=\frac1{2(b-1)}-\frac1{2b}
  \prod_{j\ge0}\left(1-b^{-2^j}\right).}                 \label{eq:kappa-b-product}
\end{equation}
For \(b=2\), this is the Prouhet--Thue--Morse constant
\begin{equation}
  \kappa_2=0.0110100110010110\ldots{}_2
  =\frac12-\frac14\prod_{j\ge0}\left(1-2^{-2^j}\right).  \label{eq:PTM-constant}
\end{equation}
The digit expansion and the lacunary product are therefore two evaluations of the same
Mahler function.

\subsection{The Woods--Robbins product}

A different global fingerprint uses the signs as exponents rather than coefficients.
The classical Woods--Robbins identity is
\begin{equation}
  \boxed{
  \prod_{n=0}^{\infty}
  \left(\frac{2n+1}{2n+2}\right)^{\eps_n}=\frac1{\sqrt2}.} \label{eq:Woods-Robbins}
\end{equation}
Convergence follows from the bounded prefix sums applied to the logarithms.  The evaluation
is obtained by separating even and odd indices, applying
\(\eps_{2n}=\eps_n\), \(\eps_{2n+1}=-\eps_n\), and simplifying the resulting functional
identity; see Woods and Robbins and the systematic extensions in
\cite{woods1978,robbins1979,alloucheriasatshallit2019}.  Many related products reduce to
gamma values, but \eqref{eq:Woods-Robbins} remains the cleanest example.

\subsection{A useful product functional equation}

For parameters for which the product converges, define
\begin{equation}
  \mathcal P(a,b)=\prod_{n\ge0}\left(\frac{n+a}{n+b}\right)^{\eps_n}.          \label{eq:parameter-product}
\end{equation}
Splitting even and odd indices gives the exact renormalization law
\begin{equation}
  \boxed{
  \mathcal P(a,b)=
  \mathcal P\!\left(\frac a2,\frac b2\right)
  \mathcal P\!\left(\frac{b+1}{2},\frac{a+1}{2}\right).}   \label{eq:parameter-product-functional}
\end{equation}
This identity is the multiplicative analogue of the Mahler equation.  Special parameter
cycles are the mechanism behind closed evaluations such as
\eqref{eq:Woods-Robbins}.

\section{Rarefied sums and Newman's phenomenon}
\label{sec:rarefied}

The bounded ordinary prefix discrepancy does not imply cancellation after restricting to an
arithmetic progression.  The most famous counterexample is the progression of multiples of
three, where one sign has a persistent excess.

\subsection{A root-of-unity filter at dyadic endpoints}

For \(q\ge2\), \(0\le r<q\), and \(m\ge0\), define
\begin{equation}
  A_{q,r}(m)=\sum_{\substack{0\le n<2^m\\n\equiv r\pmod q}}\eps_n.             \label{eq:rarefied-def}
\end{equation}
Let \(\zeta_q=\e^{2\pi\ii/q}\).  The usual root-of-unity filter and
\eqref{eq:finite-lacunary-product} give the exact finite formula
\begin{equation}
  \boxed{
  A_{q,r}(m)=\frac1q\sum_{\ell=0}^{q-1}\zeta_q^{-r\ell}
   \prod_{j=0}^{m-1}\left(1-\zeta_q^{\ell2^j}\right).}      \label{eq:rarefied-root-filter}
\end{equation}
For \(m\ge1\), the \(\ell=0\) term vanishes.  Formula
\eqref{eq:rarefied-root-filter} reduces the problem to the orbit of multiplication by two
modulo \(q\); for odd \(q\), that orbit is periodic.

\subsection{Complete evaluation modulo three}

Let \(\omega=\e^{2\pi\ii/3}\).  The powers of two alternate modulo three, and
\begin{equation}
  (1-\omega)(1-\omega^2)=3.                                \label{eq:omega-pair}
\end{equation}
This makes the rarefied sums completely explicit.

\begin{theorem}[Dyadic residue-class sums modulo three]
\label{thm:mod3-rarefied}
For \(u\ge1\),
\begin{equation}
  \boxed{
  \bigl(A_{3,0}(2u),A_{3,1}(2u),A_{3,2}(2u)\bigr)
   =\bigl(2\cdot3^{u-1},-3^{u-1},-3^{u-1}\bigr).}          \label{eq:mod3-even-level}
\end{equation}
For \(u\ge0\),
\begin{equation}
  \boxed{
  \bigl(A_{3,0}(2u+1),A_{3,1}(2u+1),A_{3,2}(2u+1)\bigr)
   =\bigl(3^u,-3^u,0\bigr).}                               \label{eq:mod3-odd-level}
\end{equation}
\end{theorem}

\begin{proof}
For \(m=2u\), the factors in \(P_m(\omega)\) occur in \(u\) pairs, each equal to three by
\eqref{eq:omega-pair}; hence
\(P_{2u}(\omega)=P_{2u}(\omega^2)=3^u\).  Substitute these values into
\eqref{eq:rarefied-root-filter} and use \(1+\omega+\omega^2=0\).

For \(m=2u+1\), one extra factor remains:
\[
  P_{2u+1}(\omega)=3^u(1-\omega),
  \qquad
  P_{2u+1}(\omega^2)=3^u(1-\omega^2).
\]
The three filtered combinations are then \(3^u,-3^u,0\).
\end{proof}

Along dyadic endpoints, the excess on multiples of three therefore grows with exponent
\begin{equation}
  \alpha=\log_4 3:\qquad
  A_{3,0}(m)\asymp (2^m)^\alpha.                            \label{eq:Newman-exponent}
\end{equation}
More precisely, the constant is \(2/3\) when \(m\) is even and \(1/\sqrt3\) when \(m\) is
odd after factoring out \((2^m)^\alpha\).

\subsection{The global Newman theorem}

For an arbitrary endpoint, put
\begin{equation}
  N_3(X)=\sum_{\substack{0\le n<X\\3\mid n}}\eps_n.        \label{eq:Newman-sum}
\end{equation}
Newman's theorem states that
\begin{equation}
  \boxed{N_3(X)>0\qquad(X\ge1).}                            \label{eq:Newman-positivity}
\end{equation}
Thus among multiples of three below every positive cutoff, those with even binary weight
outnumber those with odd binary weight.  This is not a generic balance phenomenon: it is a
rigid resonance between multiplication by two modulo three and the sign recursion.  Later
work determined sharp bounds of order \(X^{\log_4 3}\), periodic fluctuations, and
extensions to generalized digit characters and other progressions
\cite{newman1969,drmotaskalba2000,drmotastoll2008}.  The elementary dyadic formulas
\eqref{eq:mod3-even-level}--\eqref{eq:mod3-odd-level} exhibit the exponent and its two-scale
oscillation before any deeper estimate is needed.

\section{Combinatorics and dynamics of the infinite word}
\label{sec:words}

The arithmetic sign sequence and the binary word carry the same information, but the word
language exposes repetition avoidance and dynamical properties that do not look like
formulae for individual terms.

\subsection{Uniform recurrence and aperiodicity}

The substitution matrix of \(\mu(0)=01\), \(\mu(1)=10\) has every entry positive.  Hence the
substitution is primitive.  Standard primitive-substitution theory implies that the fixed
point is uniformly recurrent: every finite factor reappears with bounded gaps, and the
associated shift orbit closure is minimal and uniquely ergodic
\cite{queffelec2010}.  In particular, every factor has a well-defined frequency; the two
letters each have frequency \(1/2\).

The word is not eventually periodic.  Here is an elementary proof.  Suppose an eventual
period \(p\) existed.  Removing powers of two from \(p\) using
\(\eps_{2n}=\eps_n\) would give an odd eventual period \(q\).  Applying that period at even
indices gives, for all sufficiently large \(n\),
\[
  \eps_n=\eps_{2n}=\eps_{2n+q}
  =-\eps_{n+(q-1)/2}.
\]
Applying this relation twice shows that \(q-1\) is also an eventual period.  Since consecutive
integers have greatest common divisor one, period one would follow, making the tail constant;
this contradicts \(\eps_{2n+1}=-\eps_n\).

\subsection{Complement and reversal closure}

Complementing letters commutes with the substitution up to exchanging its two seed letters.
Therefore, whenever a finite word occurs, its bitwise complement also occurs.  The dyadic
reflection formula \eqref{eq:block-reflection} similarly shows that the reversal of every
prefix \(\mu^m(0)\) is either that prefix or its complement.  Passing to factors proves that
the full factor language is closed under both complement and reversal.

The exact balance formula \eqref{eq:exact-prefix} is the one-letter shadow of this symmetry.
For longer factors, unique ergodicity and complement invariance imply that complementary
factors have equal frequencies.

\subsection{Overlap avoidance and critical exponent}

An \emph{overlap} is a word of the form \(axaxa\), where \(a\) is one letter and \(x\) is
possibly empty.  Thue proved that the fixed point of \eqref{eq:TM-morphism} is overlap-free
\cite{thue1906}.  Consequently it contains no cube \(www\) with \(w\ne\varnothing\), and no
factor of exponent strictly larger than two.  On the other hand it contains squares of
arbitrarily large length, obtained by iterating the uniform morphism on any square factor.
Hence its critical exponent is exactly
\begin{equation}
  \boxed{2.}                                                \label{eq:critical-exponent}
\end{equation}
The overlap-free theorem is much stronger than the visible absence of \(000\) and \(111\).
It was one of the original reasons for Thue's study of the word and remains central in
combinatorics on words.

\subsection{Linear factor complexity}

Let \(p(n)\) be the number of distinct length-\(n\) factors.  Uniform morphicity gives
\(p(n)=\bigO(n)\), while aperiodicity and the Morse--Hedlund theorem give \(p(n)\ge n+1\).
Thus the word has linear, but nonminimal, factor complexity.  Exact piecewise formulas and a
complete description of square factors are known; see the survey literature in
\cite{alloucheshallit1999,alloucheshallit2003}.  This is another useful contrast: the word is
far from periodic, yet its local language grows only linearly rather than exponentially.

\section{Hankel determinants, Mahler values, and Diophantine consequences}
\label{sec:hankel}

\subsection{Hankel determinants built from the signs}

Define
\begin{equation}
  \Delta_n=\det\bigl(\eps_{i+j}\bigr)_{0\le i,j<n}.         \label{eq:Hankel-definition}
\end{equation}
The first values are
\begin{equation}
  1,-2,4,8,-16,-32,-64,128,-256,-1536,\ldots .             \label{eq:Hankel-first}
\end{equation}
Unlike the engineered Hessenberg determinant in \Cref{thm:Hessenberg-determinant}, these
matrices use no ruler weights and have no evident triangular reduction.

\begin{theorem}[Nonvanishing Hankel determinants; literature theorem]
\label{thm:Hankel-nonzero}
For every \(n\ge1\),
\begin{equation}
  \boxed{\Delta_n\ne0.}                                    \label{eq:Hankel-nonzero}
\end{equation}
\end{theorem}

This was proved by Allouche, Peyri\`ere, Wen, and Wen through a finite system of determinant
recurrences; later proofs include a short reduction modulo two and a combinatorial argument
\cite{allouchepeyrierewenwen1998,han2014,bugeaudehan2014}.  Nonvanishing means that the
coefficient series \(E(z)\) is normal for diagonal Pad\'e approximation at every order: no
unexpected rank defect destroys the corresponding rational approximant.  This determinant
rigidity is one of the main bridges from automatic sequences to Diophantine approximation.

\subsection{Transcendence of the function and of algebraic values}

The natural boundary proved in \Cref{sec:generating} immediately implies that \(E(z)\) is
transcendental over \(\C(z)\): an algebraic function has only finitely many branch points and
poles on the Riemann sphere and cannot have a natural boundary.  Mahler's method gives the
much stronger arithmetic statement.

\begin{theorem}[Mahler; literature theorem]
\label{thm:Mahler-values}
If \(\alpha\) is algebraic and \(0<|\alpha|<1\), then
\begin{equation}
  \boxed{E(\alpha)\text{ is transcendental}.}              \label{eq:Mahler-value}
\end{equation}
Consequently \(T(\alpha)\) is transcendental as well.
\end{theorem}

The hypotheses are particularly clean here: the forward orbit
\(\alpha,\alpha^2,\alpha^4,\ldots\) remains in the open unit disk and never meets the
singularity of the rational coefficient \(1-z\) in the Mahler equation
\eqref{eq:E-Mahler}.  The original result is one of the foundational examples of Mahler's
method \cite{mahler1929}.

For every integer \(b\ge2\), put
\begin{equation}
  \xi_b=T(1/b)=\sum_{n\ge0}\frac{\tau(n)}{b^n}=b\kappa_b.  \label{eq:xi-b}
\end{equation}
The theorem shows that \(\xi_b\) is transcendental.  The quality of its rational
approximations is known exactly.

\begin{theorem}[Bugeaud; literature theorem]
\label{thm:irrationality-exponent}
For every integer \(b\ge2\), the irrationality exponent of \(\xi_b\) is
\begin{equation}
  \boxed{\mu(\xi_b)=2.}                                    \label{eq:irrationality-exponent}
\end{equation}
\end{theorem}

Here \(\mu(\xi)\) is the supremum of the real numbers \(\mu\) for which
\(|\xi-p/q|<q^{-\mu}\) has infinitely many rational solutions.  The value two is the
smallest possible for an irrational real number, by Dirichlet's theorem.  Thus the
Thue--Morse--Mahler numbers are transcendental but not exceptionally well approximable.
Bugeaud's proof uses Pad\'e approximants and the nonvanishing Hankel structure
\cite{bugeaud2011}; later work placed this in a broader determinant framework
\cite{bugeaudhanwenyao2015}.

\section{The base-\texorpdfstring{\(B\)}{B} root-of-unity family}
\label{sec:baseB}

The binary sequence is the first member of a general digit-character construction.  Let
\(B\ge2\), write \(s_B(n)\) for the sum of the base-\(B\) digits of \(n\), and fix
\(\zeta\in\C\).  Define
\begin{equation}
  u_n=\zeta^{s_B(n)}.                                       \label{eq:baseB-character}
\end{equation}
For \(0\le n<B^m\), choosing the \(m\) digits independently gives the finite product
\begin{equation}
  \boxed{
  \sum_{n=0}^{B^m-1}u_nz^n
   =\prod_{j=0}^{m-1}
     \left(1+\zeta z^{B^j}+\zeta^2z^{2B^j}+\cdots+
           \zeta^{B-1}z^{(B-1)B^j}\right).}                \label{eq:baseB-product}
\end{equation}
It also gives block multiplicativity
\begin{equation}
  u_{B^mq+r}=u_qu_r\qquad(0\le r<B^m).                     \label{eq:baseB-block}
\end{equation}

Suppose now that \(\zeta\) is a nontrivial \(B\)-th root of unity.  At \(z=\e^t\), each
factor in \eqref{eq:baseB-product} has a simple zero at \(t=0\), because
\begin{equation}
  \sum_{d=0}^{B-1}\zeta^d=0,
  \qquad
  \sum_{d=0}^{B-1}d\zeta^d=-\frac{B}{1-\zeta}\ne0.         \label{eq:root-unity-derivative}
\end{equation}
Thus the generalized character annihilates powers below degree \(m\).

\begin{theorem}[Base-\(B\) Prouhet character]
\label{thm:baseB-Prouhet}
If \(\zeta^B=1\) and \(\zeta\ne1\), then
\begin{align}
  \sum_{n=0}^{B^m-1}\zeta^{s_B(n)}n^r&=0 &&(0\le r<m),     \label{eq:baseB-Prouhet-zero}\\
  \sum_{n=0}^{B^m-1}\zeta^{s_B(n)}n^m
   &=m!\left(-\frac{B}{1-\zeta}\right)^m
       B^{\binom m2}.                                      \label{eq:baseB-Prouhet-boundary}
\end{align}
\end{theorem}

\begin{proof}
Set \(z=\e^t\) in \eqref{eq:baseB-product}.  The derivative at zero of the \(j\)-th factor
is \(-B^{j+1}/(1-\zeta)\).  Their product is the coefficient of \(t^m\), namely
\((-B/(1-\zeta))^mB^{\binom m2}\).  Compare exponential generating coefficients.
\end{proof}

For \(B=2\) and \(\zeta=-1\), one recovers \(u_n=\eps_n\), the lacunary factors
\(1-z^{2^j}\), and the exact moment \eqref{eq:first-surviving-moment}.  This generalization
shows which parts of the atlas arise from digit independence and root-of-unity cancellation,
and which depend specifically on binary parity, Lucas's theorem modulo two, or two-adic
valuations.

\section{Formula selection and a Lean formalization roadmap}
\label{sec:formalization}

\subsection{Which representation is useful where?}

No single formula dominates all others.  The right interface depends on the data already
present in a proof or computation.

\renewcommand{\arraystretch}{1.16}
\begin{longtable}{@{}p{0.24\textwidth}p{0.31\textwidth}p{0.37\textwidth}@{}}
\toprule
Representation & Required input & Best use \\
\midrule
\endfirsthead
\toprule
Representation & Required input & Best use \\
\midrule
\endhead
Binary popcount & machine word or bit list & Fast evaluation; executable normalization. \\
Even--odd recursion & quotient and low bit & Induction; automata; streaming generation. \\
Dyadic block law & high and low bit blocks & Kernel arguments; divide-and-conquer proofs. \\
XOR character & bitwise XOR & Boolean harmonic analysis; carry-free sums. \\
Factorial valuation & \(n!\) or floor sums & Number-theoretic contexts already using valuations. \\
Central binomial valuation & \(\binom{2n}{n}\) & A one-integer formula; Kummer interfaces. \\
Carry cocycle & addition and a binomial valuation & Translating XOR identities to ordinary addition. \\
Pascal support & parity of row entries & Lucas-theorem and combinatorial arguments. \\
Modulo-three odd count & \(\nu(n)\bmod3\) & Log-free version of the Pascal-row formula. \\
Finite lacunary product & polynomial or formal series & Prouhet, cyclotomic, Fourier, and coefficient proofs. \\
Ruler convolution & earlier signs and \(v_2(k)\) & Digit-free recurrence; determinant extraction. \\
Finite differences & arbitrary test function & Cancellation beyond monomials; positivity and bounds. \\
Walsh transform & Boolean vectors & One-frequency spectral normal form. \\
Ordinary DFT & roots of unity & Cyclic spectra and rarefied sums. \\
Riesz product & Fourier intensity & Spectral measures and autocorrelation. \\
Dirichlet/Mellin transform & analytic continuation & Global analytic fingerprints and special values. \\
\bottomrule
\end{longtable}
\renewcommand{\arraystretch}{1}

\subsection{Existing formalized spine}

The repository's proved development already contains the core chain
\begin{equation}
  \text{binary recursion}
  \Longrightarrow \text{block product}
  \Longrightarrow \text{Prouhet cancellation}
  \Longrightarrow \text{dyadic Fabius arithmetic}.         \label{eq:formal-spine}
\end{equation}
The principal files relevant to this article include
\begin{itemize}
\item \lean{ThueMorseGenerating.lean}: finite block polynomial and lacunary product;
\item \lean{ThueMorseExponential.lean}: exponential form, order of vanishing, and moments;
\item \lean{ThueMorsePrefix.lean}: iterated-prefix identities;
\item \lean{ThueMorseApproximation.lean} and
      \lean{PaperKFoldThueMorse.lean}: convergence of discrete prefix constructions;
\item \lean{ThueMorseBinomialLog.lean}: the Pascal-row logarithm;
\item \lean{DyadicClosedForm.lean}, \lean{DyadicCorrectness.lean}, and the
      \lean{Fabius*QBinomial*.lean} files: exact dyadic evaluation and its
      \(q\)-binomial forms.
\end{itemize}
The main exposition also names \lean{binaryWeight}, \lean{thueMorseSign}, and the
primitive-recursive bit evaluator \lean{tmBitPR} as the foundational term-level objects.

\subsection{Suggested dependency order for the new formulae}

A practical formalization should avoid beginning with the analytic or determinant results.
The following order maximizes reuse.

\begin{enumerate}
\item \textbf{Digit and bit API.}  Prove floor extraction, XOR parity, block reflection,
      and the ruler transition from existing binary-weight lemmas.
\item \textbf{Valuations.}  Establish the base-two Legendre identity and then central
      binomial, Catalan, Kummer, and multinomial corollaries.
\item \textbf{Pascal support.}  Specialize Lucas modulo two, identify submasks, count odd
      entries, and prove the modulo-three and \(r\)-nomial formulas.
\item \textbf{Boolean incidence algebra.}  Formalize the submask zeta transform and the
      explicit Thue--Morse M\"obius inverse.
\item \textbf{Euler-transform layer.}  Starting from the formal product, derive the ruler
      logarithm, convolution recurrence, Bell formula, and Hessenberg determinant.
\item \textbf{General finite differences.}  Lift polynomial Prouhet cancellation to
      arbitrary commuting translations and derive sparse Pascal-row cancellation.
\item \textbf{Finite harmonic analysis.}  Evaluate the block polynomial on roots of unity,
      prove the DFT support and Walsh delta theorem, then derive the modulo-three rarefied
      formulas.
\item \textbf{Analysis.}  Add the Mellin continuation and parameterized infinite products
      only after the formal power-series and finite-sum layers are stable.
\end{enumerate}

Natural module names would be
\begin{center}
\small
\begin{tabularx}{0.95\textwidth}{@{}lX@{}}
\toprule
Proposed module & Principal declarations \\
\midrule
\lean{ThueMorseValuation.lean} &
Legendre digit sum, central-binomial valuation, Catalan parity, carry cocycle.\\
\lean{ThueMorsePascalSupport.lean} &
Lucas submask criterion, Glaisher count, sparse columns, modulo-three formula.\\
\lean{ThueMorseBooleanMobius.lean} &
Submask incidence matrix and explicit inverse.\\
\lean{ThueMorseEulerTransform.lean} &
Ruler logarithm, convolution recurrence, Bell and determinant formulae.\\
\lean{ThueMorseSparseProuhet.lean} &
Arbitrary-support product and exact boundary moment.\\
\lean{ThueMorseFiniteFourier.lean} &
DFT product, even-frequency vanishing, Walsh transform.\\
\lean{ThueMorseRarefied.lean} &
Root-of-unity filter and exact modulo-three dyadic sums.\\
\lean{ThueMorseDirichlet.lean} &
Mellin representation, superpolynomial endpoint decay, entire continuation.\\
\bottomrule
\end{tabularx}
\end{center}

\subsection{Executable checks}

For evaluation, popcount parity remains the correct implementation.  The other formulae are
valuable as proof interfaces and independent regression oracles.

\begin{lstlisting}[style=pseudo,caption={Cross-checking independent formula families}]
eps(n):
    return +1 if popcount(n) is even else -1

for n from 0 to limit:
    e = eps(n)
    assert e == (-1)^(n - v2(factorial(n)))
    assert e == (-1)^v2(binomial(2*n, n))
    odd = count(k in 0..n where binomial(n,k) is odd)
    assert odd == 2^popcount(n)
    assert e == 3 - 2*(odd mod 3)

for m from 0 to maxLevel:
    P = sum(n in 0..2^m-1, eps(n)*z^n)
    assert P == product(j in 0..m-1, 1-z^(2^j))
    for r from 0 to m-1:
        assert sum(n in 0..2^m-1, eps(n)*n^r) == 0
    assert sum(n in 0..2^m-1, eps(n)*n^m)
           == (-1)^m * factorial(m) * 2^(m*(m-1)/2)
\end{lstlisting}

For fixed-width machine integers, the optimal evaluator is a parity-popcount instruction or a
sequence of XOR folds.  The recursion is optimal for an unbounded inductive natural.  The
central-binomial and Pascal formulas should not be used to compute large terms literally;
their role is to expose mathematical structure and to connect with existing arithmetic data.

\section{Worked examples and consistency checks}
\label{sec:examples}

\subsection{One index through ten formulae}

Take \(n=13=(1101)_2\).  Its weight is three, so \(\tau(13)=1\) and \(\eps_{13}=-1\).
The same value emerges as follows.

\begin{enumerate}
\item \textbf{Digit product:}
\((1-2)(1-0)(1-2)(1-2)=-1\).
\item \textbf{Floor sum:}
\(\lfloor13/2\rfloor+\lfloor13/4\rfloor+\lfloor13/8\rfloor=6+3+1=10\), so
\((-1)^{13-10}=-1\).
\item \textbf{Factorial valuation:}
\(\vtwo(13!)=10\), again giving \((-1)^3\).
\item \textbf{Central binomial:}
\(\vtwo\binom{26}{13}=\wt(13)=3\).
\item \textbf{Catalan:}
\(\vtwo(C_{13})=\wt(14)-1=2\) and \(\vtwo(14)=1\), whose sum is three.
\item \textbf{Sparse Pascal columns:}
\(\binom{13}{1},\binom{13}{4},\binom{13}{8}\) are odd, while
\(\binom{13}{2}\) is even; the parity sum is three.
\item \textbf{Pascal support:}
row 13 has \(2^3=8\) odd entries, and \(8\bmod3=2\), so
\(3-2(2)=-1\).
\item \textbf{Submask product:}
\[
 \sum_{k\submask13}\eps_kz^k=(1-z)(1-z^4)(1-z^8),
\]
whose coefficient at \(z^{13}\) is \(-1\).
\item \textbf{Ruler transition:}
starting at \(\eps_0=1\), repeatedly apply
\(\eps_{j+1}=-(-1)^{\vtwo(j+1)}\eps_j\) through \(j=12\).
\item \textbf{Hessenberg determinant:}
with \(a_j=2^{\vtwo(j)+1}-1\),
\(\det H_{13}=(-1)^{13}13!\eps_{13}=13!\).
\end{enumerate}
The routes use radically different intermediate objects, but every one reduces to the same
binary parity character.

\subsection{A complete four-bit block}

At \(m=4\),
\begin{align}
  P_4(z)
  &=1-z-z^2+z^3-z^4+z^5+z^6-z^7                     \notag\\
  &\quad{}-z^8+z^9+z^{10}-z^{11}+z^{12}-z^{13}-z^{14}+z^{15}                  \notag\\
  &=(1-z)(1-z^2)(1-z^4)(1-z^8).                          \label{eq:P4-example}
\end{align}
The factor \((1-z)^4\) divides this polynomial, so its first four ordinary derivatives at
one vanish through order three.  Equivalently,
\begin{equation}
  \sum_{n<16}\eps_nn^r=0\quad(0\le r<4),
  \qquad
  \sum_{n<16}\eps_nn^4=4!\,2^6=1536.                       \label{eq:P4-moments}
\end{equation}
The sign is positive because \((-1)^4=1\).

The residue-class sums modulo three are, from \eqref{eq:mod3-even-level},
\begin{equation}
  \left(
  \sum_{\substack{n<16\\n\equiv0(3)}}\eps_n,
  \sum_{\substack{n<16\\n\equiv1(3)}}\eps_n,
  \sum_{\substack{n<16\\n\equiv2(3)}}\eps_n
  \right)=(6,-3,-3).                                      \label{eq:P4-rarefied}
\end{equation}
Their sum is zero, as required by the ordinary block balance, but the distribution among
residue classes is highly nonuniform.

\subsection{A carry-corrected addition}

Take \(a=13=(1101)_2\) and \(b=7=(0111)_2\).  Then
\(\eps_{13}=\eps_7=-1\), while \(a+b=20=(10100)_2\) has sign \(+1\).  Kummer's formula gives
\begin{equation}
  \vtwo\binom{20}{13}=\wt(13)+\wt(7)-\wt(20)=3+3-2=4.      \label{eq:carry-example}
\end{equation}
The correction factor \((-1)^4=1\) precisely restores multiplicativity.  This example also
illustrates why a yes/no statement about the presence of carries is insufficient: their
parity, with propagation counted, is the invariant required by the sign.

\section{Concluding synthesis}
\label{sec:conclusion}

The Thue--Morse sequence is often introduced by the sentence ``take the parity of the number
of one-bits.''  Correct as that is, it conceals the main structural fact.  The same parity
character is simultaneously
\begin{itemize}
\item a two-state automatic sequence and a character of the Boolean XOR group;
\item the parity of a factorial, central-binomial, Catalan, or carry valuation;
\item the M\"obius kernel of the Boolean support of Pascal's triangle;
\item the coefficient sequence of a lacunary Mahler product;
\item an inclusion--exclusion kernel for the dyadic differences
      \(1,2,4,\ldots\);
\item Prouhet's exact annihilator of low-degree polynomials;
\item one Walsh frequency but a singular continuous ordinary Fourier object;
\item the sign system in the global Fabius translate expansion and in exact dyadic
      cancellation;
\item the source of rarefied arithmetic-progression biases, Hankel rigidity, and
      transcendental Mahler values.
\end{itemize}
These are not unrelated coincidences.  Binary digit independence gives the finite product;
Lucas turns that independence into Pascal support; Legendre and Kummer turn it into
valuations and carries; evaluating the product at \(1\), roots of unity, exponentials, or
translation operators yields Prouhet, Fourier, rarefied, and finite-difference identities.
The binomial--logarithm formula sits exactly where the Pascal and digit branches meet.

The deepest reusable principle is therefore not any one boxed expression, but the following
commuting diagram of transformations:
\begin{equation}
\begin{gathered}
  \text{binary digits}
  \longleftrightarrow \text{Boolean character}
  \longleftrightarrow \text{Lucas support},\\[0.3em]
  \text{binary digits}
  \longleftrightarrow \text{Legendre--Kummer valuations},\\[0.3em]
  \text{Boolean character}
  \longleftrightarrow \prod_j(1-z^{2^j})
  \longrightarrow
  \begin{cases}
    \text{finite differences and Prouhet moments},\\
    \text{ordinary Fourier and Riesz products},\\
    \text{Mahler, Mellin, and Dirichlet transforms},\\
    \text{root-of-unity filters and rarefied sums}.
  \end{cases}
\end{gathered}                                              \label{eq:final-diagram}
\end{equation}
For the Fabius project, this diagram provides several independent proof interfaces to the
same arithmetic layer.  For future formalization, it also suggests a clean order: digits,
valuations and Lucas support first; formal products and finite differences next; spectral and
analytic consequences last.

\appendix

\section{The first 64 terms}
\label{app:first-64}

The following table gives \(\tau(n)\).  Replacing each zero by \(+1\) and each one by
\(-1\) gives \(\eps_n\).

\begin{center}
\small
\begin{tabular}{r|rrrrrrrr}
\toprule
start & \(+0\) & \(+1\) & \(+2\) & \(+3\) & \(+4\) & \(+5\) & \(+6\) & \(+7\)\\
\midrule
0  &0&1&1&0&1&0&0&1\\
8  &1&0&0&1&0&1&1&0\\
16 &1&0&0&1&0&1&1&0\\
24 &0&1&1&0&1&0&0&1\\
32 &1&0&0&1&0&1&1&0\\
40 &0&1&1&0&1&0&0&1\\
48 &0&1&1&0&1&0&0&1\\
56 &1&0&0&1&0&1&1&0\\
\bottomrule
\end{tabular}
\end{center}

The first 32-bit block is followed by its complement.  Inside each 16-bit and 8-bit block,
the same copy/complement rule recurses.  This nested visual structure is precisely
\eqref{eq:tensor-block}.

\section{A compact identity map}
\label{app:identity-map}

For quick reference, the main term-level identities are
\begin{align*}
  \eps_n
  &=(-1)^{\wt(n)}
   =\prod_{j\ge0}(1-2b_j(n))\\
  &=(-1)^{n-\sum_{j\ge1}\lfloor n/2^j\rfloor}
   =(-1)^{n-\vtwo(n!)}\\
  &=(-1)^{\vtwo\binom{2n}{n}}
   =(-1)^{\sum_{j\ge0}\binom{n}{2^j}}\\
  &=\prod_{j\ge0}\cos\!\left(\pi\binom{n}{2^j}\right)
   =\prod_{j\ge0}\cos\!\left(\pi\left\lfloor\frac n{2^j}\right\rfloor\right)\\
  &=\prod_{j\ge0}\sgn\!\sin\!\left(\frac{\pi(n+1/2)}{2^j}\right)\\
  &=(-1)^{\log_2\nu(n)}
   =3-2\res{\nu(n)}3\\
  &=[z^n]\prod_{j\ge0}(1-z^{2^j})
   =\frac{(-1)^n}{n!}\det H_n.
\end{align*}
The principal block identities are
\begin{align*}
  \sum_{n<2^m}\eps_nz^n
   &=\prod_{j<m}(1-z^{2^j}),\\
  \sum_{n<2^m}\eps_n\e^{nt}
   &=\prod_{j<m}(1-\e^{2^jt}),\\
  \sum_{n<2^m}\eps_nf(x+nh)
   &=(-1)^m\Delta_h\Delta_{2h}\cdots\Delta_{2^{m-1}h}f(x),\\
  \sum_{n<2^m}\eps_nn^r&=0\quad(r<m),\\
  \sum_{n<2^m}\eps_nn^m&=(-1)^mm!2^{\binom m2},\\
  \widehat\eps_m(k)&=\prod_{j<m}
     (1-\e^{-2\pi\ii k2^j/2^m}),\\
  2^{-m}|P_m(\e^{2\pi\ii x})|^2
   &=\prod_{j<m}(1-\cos(2\pi2^jx)).
\end{align*}
Each line is obtained from the same finite polynomial by coefficient extraction,
substitution, differentiation, or a change of harmonic-analysis basis.

\section{Consolidation provenance}
\label{app:provenance}

The four source drafts had substantial duplicated text but also distinct mathematical
contributions.  The consolidation used the following division of labor.

\begin{itemize}
\item \texttt{Further\_Formulae\_for\_the\_Thue\_Morse\_Sequence-1.tex} supplied the strongest
      Boolean M\"obius-inversion, ruler-convolution, Bell-polynomial, Hessenberg-determinant,
      finite-Fourier, and Dirichlet-transform material.
\item \texttt{...-2.tex} supplied the master-catalogue organization, the carry-corrected
      addition law, mixed finite differences, all-order moment formulas, computational
      comparison, and the base-\(B\) discovery pattern.
\item \texttt{...-3.tex} supplied the \(r\)-nomial odd-coefficient family, sparse Prouhet
      cancellation on arbitrary Pascal support, the characteristic-two algebraic series,
      and a particularly clean separation between elementary proofs and literature results.
\item \texttt{...-4.tex} supplied the logarithm-free modulo-three Pascal formula, odious and
      evil enumerators, the integer algebraic lift, autocorrelation and Riesz-product
      organization, and several exact integral and product forms.
\end{itemize}
The present article adopts one notation, one theorem numbering system, one bibliography, and
one proof-status convention.  Repeated elementary arguments were retained only once, while
stronger versions replaced weaker duplicates.  Sections on rarefied sums, combinatorics on
words, Hankel determinants, natural boundaries, Mahler values, and irrationality exponents
were added during consolidation.

\begin{thebibliography}{99}

\bibitem{allouchepeyrierewenwen1998}
J.-P. Allouche, J. Peyri\`ere, Z.-X. Wen, and Z.-Y. Wen,
\emph{Hankel determinants of the Thue--Morse sequence},
Annales de l'Institut Fourier \textbf{48} (1998), 1--27.

\bibitem{alloucheshallit1999}
J.-P. Allouche and J. Shallit,
\emph{The ubiquitous Prouhet--Thue--Morse sequence},
in C. Ding, T. Helleseth, and H. Niederreiter (eds.),
\emph{Sequences and Their Applications}, Springer, London, 1999, 1--16.

\bibitem{alloucheshallit2003}
J.-P. Allouche and J. Shallit,
\emph{Automatic Sequences: Theory, Applications, Generalizations},
Cambridge University Press, Cambridge, 2003.

\bibitem{alloucheriasatshallit2019}
J.-P. Allouche, S. Riasat, and J. Shallit,
\emph{More infinite products: Thue--Morse and the gamma function},
The Ramanujan Journal \textbf{49} (2019), 115--128;
\url{https://arxiv.org/abs/1709.03398}.

\bibitem{baakegrimm2008}
M. Baake and U. Grimm,
\emph{The singular continuous diffraction measure of the Thue--Morse chain},
Journal of Physics A: Mathematical and Theoretical \textbf{41} (2008), 422001;
\url{https://arxiv.org/abs/0809.0580}.

\bibitem{bugeaud2011}
Y. Bugeaud,
\emph{On the rational approximation to the Thue--Morse--Mahler numbers},
Annales de l'Institut Fourier \textbf{61} (2011), no.~5, 2065--2076;
\url{https://doi.org/10.5802/aif.2666}.

\bibitem{bugeaudehan2014}
Y. Bugeaud and G.-N. Han,
\emph{A combinatorial proof of the non-vanishing of Hankel determinants of the
Thue--Morse sequence},
arXiv:1406.1587 (2014);
\url{https://arxiv.org/abs/1406.1587}.

\bibitem{bugeaudhanwenyao2015}
Y. Bugeaud, G.-N. Han, Z.-Y. Wen, and J.-Y. Yao,
\emph{Hankel determinants, Pad\'e approximations, and irrationality exponents},
arXiv:1503.02797 (2015);
\url{https://arxiv.org/abs/1503.02797}.

\bibitem{christol1980}
G. Christol, T. Kamae, M. Mend\`es France, and G. Rauzy,
\emph{Suites alg\'ebriques, automates et substitutions},
Bulletin de la Soci\'et\'e Math\'ematique de France \textbf{108} (1980), 401--419.

\bibitem{drmotaskalba2000}
M. Drmota and M. Ska\l ba,
\emph{Rarified sums of the Thue--Morse sequence},
Transactions of the American Mathematical Society \textbf{352} (2000), 609--642.

\bibitem{drmotastoll2008}
M. Drmota and T. Stoll,
\emph{Newman's phenomenon for generalized Thue--Morse sequences},
Discrete Mathematics \textbf{308} (2008), 1191--1208.

\bibitem{fine1947}
N. J. Fine,
\emph{Binomial coefficients modulo a prime},
The American Mathematical Monthly \textbf{54} (1947), 589--592.

\bibitem{han2014}
G.-N. Han,
\emph{Hankel determinant calculus for the Thue--Morse and related sequences},
arXiv:1406.1586 (2014);
\url{https://arxiv.org/abs/1406.1586}.

\bibitem{kummer1852}
E. E. Kummer,
\emph{\"Uber die Erg\"anzungss\"atze zu den allgemeinen Reciprocit\"atsgesetzen},
Journal f\"ur die reine und angewandte Mathematik \textbf{44} (1852), 93--146.

\bibitem{mahler1929}
K. Mahler,
\emph{Arithmetische Eigenschaften der L\"osungen einer Klasse von Funktionalgleichungen},
Mathematische Annalen \textbf{101} (1929), 342--366.

\bibitem{morse1921}
M. Morse,
\emph{Recurrent geodesics on a surface of negative curvature},
Transactions of the American Mathematical Society \textbf{22} (1921), 84--100.

\bibitem{newman1969}
D. J. Newman,
\emph{On the number of binary digits in a multiple of three},
Proceedings of the American Mathematical Society \textbf{21} (1969), 719--721.

\bibitem{prouhet1851}
E. Prouhet,
\emph{M\'emoire sur quelques relations entre les puissances des nombres},
Comptes rendus de l'Acad\'emie des sciences de Paris \textbf{33} (1851), 225.

\bibitem{queffelec2010}
M. Queff\'elec,
\emph{Substitution Dynamical Systems -- Spectral Analysis},
2nd ed., Lecture Notes in Mathematics 1294, Springer, Berlin, 2010.

\bibitem{reshetnikovrepo}
V. Reshetnikov,
\emph{ProveIt: Analysis/FabiusFunction}, Lean~4 formal development,
repository snapshot consulted 26 August 2026;
\url{https://github.com/VladimirReshetnikov/ProveIt/tree/main/Analysis/FabiusFunction}.

\bibitem{reshetnikovexposition}
V. Reshetnikov,
\emph{The Fabius Function and Rvachev's Up-Function: Exact arithmetic, probability,
Fourier analysis, and corrected sharp asymptotics},
source in the \texttt{ProveIt} repository, especially the exact dyadic and
Thue--Morse chapters;
\url{https://github.com/VladimirReshetnikov/ProveIt/blob/main/Analysis/FabiusFunction/docs/Fabius_Function_and_Rvachev_Up/Fabius_Function_and_Rvachev_Up.tex}.

\bibitem{robbins1979}
D. Robbins,
\emph{Solution to elementary problem E2692},
The American Mathematical Monthly \textbf{86} (1979), 394--395.

\bibitem{thue1906}
A. Thue,
\emph{\"Uber unendliche Zeichenreihen},
Norske Videnskabers Selskabs Skrifter, I. Matematisk-naturvidenskabelig Klasse,
no.~7 (1906), 1--22.

\bibitem{woods1978}
D. R. Woods,
\emph{Elementary problem E2692},
The American Mathematical Monthly \textbf{85} (1978), 48.

\end{thebibliography}

\end{document}
